\section{A.5. The octahedron and four unordered points}

\paragraph{Conventions and source control.}
We use scan numbers, so that the manuscript's theorem is at cote~67/089.
An unordered four-point curve means a smooth proper geometrically connected
genus-zero curve together with a reduced degree-four divisor, and its
automorphisms may permute the divisor. A labelled curve has four separate
labelled sections and all its automorphisms fix those sections. These are
different moduli problems. All statements in this subsection that use an
octahedron assume that the characteristic is different from~2.

\begin{proposition}[Proved: the corrected proposition of 67/086]
\label{A:octahedron}
Let $k$ be an algebraically closed field with $\operatorname{char}k\ne2$,
let $G=\operatorname{PGL}_2(k)$, and choose $i\in k$ with $i^2=-1$.
Put
\[
 V=\{z\mapsto z,\ z\mapsto-z,\ z\mapsto1/z,\ z\mapsto-1/z\},
 \qquad S=\{0,\infty,1,-1,i,-i\}.
\]
Every Klein four-subgroup of $G$ is conjugate to $V$. Its normalizer has
order~24 and is isomorphic to $\mathfrak S_4$. The assignments
$V'\mapsto N_G(V')$ and $\Gamma\mapsto$ its normal Klein four-subgroup
give inverse bijections between Klein four-subgroups and subgroups
isomorphic to $\mathfrak S_4$, in every characteristic different from~2.
The assignment taking $V'$ to the union of the fixed points of its three
nonidentity elements is a bijection onto the projective translates of
$S$ if and only if $\operatorname{char}k\ne5$. In characteristic~5 it is
five-to-one, and the stabilizer of $S$ has order~120.
\end{proposition}

\begin{proof}
Every involution of $\operatorname{PGL}_2(k)$ is diagonalizable with two
distinct fixed points, and after taking those points to $0,\infty$ it is
$z\mapsto-z$. A commuting second involution either fixes both points,
in which case it is the same involution, or interchanges them, in which
case it has the form $z\mapsto a/z$. Replacing $z$ by $bz$ with a suitable
$b^2$ reduces this to $z\mapsto1/z$. This proves the conjugacy assertion.

The centralizer of $V$ is $V$: an element commuting with $z\mapsto-z$
is $z\mapsto az$ or $z\mapsto a/z$, and commuting also with
$z\mapsto1/z$ forces $a=\pm1$. Thus conjugation embeds
$N_G(V)/V$ into $\operatorname{Aut}(V)\simeq\mathfrak S_3$.
The transformations
\[
 r(z)=iz,\qquad h(z)=\frac{z+1}{z-1}
\]
normalize $V$ and give the two transpositions of its three nonidentity
elements. Consequently the normalizer has order $4\cdot6=24$.
To identify it, partition $S$ into the opposite pairs
$\{0,\infty\}$, $\{1,-1\}$, $\{i,-i\}$. The eight choices of one point
from each pair give four unordered opposite pairs of faces. Both $r$
and $h$ permute these four pairs, and the resulting action of the group
they generate is faithful: in the signed permutation description of the
six vertices they are the orientation-preserving octahedral permutations,
and the kernel on the four body diagonals is trivial. Explicitly, label
an opposite face pair by $(\epsilon_1,\epsilon_2,\epsilon_3)$ modulo
simultaneous sign reversal; the four sign classes are permuted faithfully
by the 24 signed coordinate permutations of determinant~1. The displayed
transformations have exactly these permutations, independently of the
characteristic. Hence $N_G(V)\simeq\mathfrak S_4$.

An embedded $\mathfrak S_4$ has its characteristic normal subgroup of
double transpositions, which is a Klein group. The embedded group is
contained in the normalizer of that subgroup and both have order~24,
so they are equal. This proves the first bijection.

Write $H=\operatorname{Stab}_G(S)$. The normalizer acts transitively on
$S$: $V$ and $h$ move $0$ to $0,\infty,1,-1$, and $r$ supplies $\pm i$.
An element of $H$ fixing $\infty$ is $z\mapsto az+b$ and permutes
$T=\{0,1,-1,i,-i\}$. Summing the five elements of $T$ before and after
the permutation gives $5b=0$. Outside characteristic~5, $b=0$; then
$aT=T$ if and only if $a\in\{1,-1,i,-i\}$. Thus $|H|=6\cdot4=24$
and $H=N_G(V)$. This argument includes characteristic~3.

If another Klein group $V'$ has fixed-point set $S$, a projective
conjugacy from $V$ to $V'$ must preserve $S$. Thus it lies in $H$.
Outside characteristic~5 this is $N_G(V)$, and $V'=V$.
In characteristic~5 we may take $i=2$, so $S=\mathbb P^1(\mathbb F_5)$.
A projective transformation preserving this set is determined by the
images of three distinct rational points; the unique transformation
with those images is defined over $\mathbb F_5$. Therefore
$H=\operatorname{PGL}_2(\mathbb F_5)$, of order $5(5^2-1)=120$.
The same conjugacy argument shows that all Klein groups with this
fixed-point set form one $H$-orbit, whose stabilizer is $N_G(V)$ of
order~24. There are exactly five. This also gives five distinct
octahedral normalizers over the same six-point set.
\end{proof}

\begin{remark}[Refuted statement and finite certificate]
\label{A:octahedron-certificate}
The characteristic-$\ne2$ bijection with bare six-point sets in
67/084 and 67/086 is false. The explicit counterexample is
$\mathbb P^1_{\mathbb F_5}$ with its six rational points, after base
extension to $\overline{\mathbb F}_5$. All five Klein groups can be
written without a search: for $b\in\mathbb F_5$ let
\[
 V_b=\left\{z,\ 2b-z,\ b+\frac1{z-b},\ b-\frac1{z-b}\right\}.
\]
They are the conjugates of $V$ by $z\mapsto z+b$. Two coincide only
if their relative nontrivial translation belongs to $N_G(V)$, which
is impossible since such a translation has order~5 and $|N_G(V)|=24$.
They all have fixed-point union $S$, and the preceding orbit proof
shows that this list is exhaustive. The program
\texttt{code/A567\_octahedral.py} enumerates all 120 scalar-normalized
invertible $2\times2$ matrices over $\mathbb F_5$, all 25 involutions,
and all 20 Klein subgroups. Exactly five have their six geometric fixed
points in this set; their five normalizers each have order~24.
The output gives all five groups as matrices. The computation is an
independent finite check of the counterexample; the preceding proof,
not a search bound, excludes every other characteristic.
In characteristic~2 the six points are not distinct and the
diagonalization step for an involution fails. No claim here extends
the octahedral formulation to characteristic~2.
\end{remark}

\begin{theorem}[Proved: the four-point groupoid of 67/089]
\label{A:fourpoint-groupoid}
Let $k$ be algebraically closed of characteristic different from~2.
For a four-element set $I$, let $V(I)$ be its canonical Klein group of
double transpositions and let $J(I)=V(I)\setminus\{1\}$.
Let $Z(J)$ be the projective line with its three labelled distinct
points indexed by $J$, and write $Z(J)^*$ for the complement of those
points. There is an equivalence between the groupoid of four-point
curves $(X,I\hookrightarrow X)$, with bijections of $I$ allowed, and
the groupoid of $(I,\zeta)$ with $\zeta\in Z(J(I))^*$.
This remains true in characteristic~5. The construction retains
$V(I)$ and does not attempt to recover it from a bare six-point set.
\end{theorem}

\begin{proof}
The double transpositions of four distinct points extend uniquely to
projective automorphisms. For example, after putting the points at
$\infty,0,1,\lambda$, one is $z\mapsto\lambda/z$; the other two
are obtained by relabelling. They form $V(I)$ because the extension
of a permutation, if it exists, is unique. None fixes one of the four
points. The quotient $q:X\to X/V(I)$ has degree~4 and three branch
points labelled by the three nontrivial stabilizers. In the standard
model it is $q(z)=z^2+z^{-2}$, with branch values $\infty,2,-2$.
These are distinct precisely under our characteristic hypothesis.
There is a unique isomorphism of the three-point quotient with
$Z(J(I))$. The image of $I$ is therefore a well-defined point
$\zeta\in Z(J(I))^*$.

Here is an inverse that keeps track of all automorphisms. Choose once
an octahedral model $X_0(I)$ with its $V(I)$-action and quotient
$Z(J(I))$. It can be chosen functorially in $I$: choose a standard
four-element set $I_0$, identify its permutation group with the
normalizer in Proposition~\ref{A:octahedron} by the action on the four
opposite face pairs, and transport the projective line along
$\operatorname{Isom}(I_0,I)$. The permutations of $I$ act on $V(I)$
by conjugation, so this transport also transports the quotient labels.

The fibre $F_\zeta=X_0(I)_\zeta$ is a $V(I)$-torsor. So is $I$.
Let $T=\operatorname{Isom}_{V(I)}(F_\zeta,I)$, another $V(I)$-torsor.
The contracted product
\[
 X=X_0(I)\mathbin{\wedge}_{V(I)}T
\]
is a projective line, with a canonical identification of its fibre at
$\zeta$ with $I$. To verify the inverse property, any two projective
lines with the prescribed faithful $V(I)$-action are equivariantly
isomorphic by the first paragraph of the proof of
Proposition~\ref{A:octahedron}. The set of equivariant isomorphisms
is a torsor under the centralizer, which that proof computed as
$V(I)$. Prescribing its restriction to the free four-point fibre
therefore prescribes exactly the torsor $T$ above. This proves
essential surjectivity and full faithfulness, not only a bijection of
isomorphism classes. Every step is natural under bijections of $I$.
In characteristic~5 the same centralizer and normalizer computations
hold, so the larger stabilizer of the unlabelled six-point set is
irrelevant.
\end{proof}

\begin{proposition}[Proved: the stack and coarse space in Lead L7]
\label{A:unordered-stack}
Put $R=\mathbb Z[1/2,i]$, and
$U=\operatorname{Spec}R[\lambda,\lambda^{-1},(1-\lambda)^{-1}]$.
The stack of genus-zero curves with an unordered finite \'etale
degree-four divisor is
\[
 \mathcal M_{0,[4],R}\simeq[U/\mathfrak S_4],
 \qquad
 \mathfrak S_4\longrightarrow\mathfrak S_3
       =\operatorname{Aut}\{0,1,\infty\}.
\]
The action factors through the displayed quotient and has kernel
$V_4$. Its coarse quotient is
\[
 \operatorname{Spec}R[j],\qquad
 j=256\frac{(\lambda^2-\lambda+1)^3}
                    {\lambda^2(1-\lambda)^2}.
\]
Consequently Lead L7 is correct as a statement about the unordered
moduli stack, together with its coarse space, but a groupoid of
unordered four-point curves is not represented by the scheme
$\mathbb A^1_R$. The stack description in fact descends to
$\mathbb Z$; adjoining $i$ and inverting~2 are needed only for the
specified octahedral presentation.
\end{proposition}

\begin{proof}
Ordering the four points is an \'etale $\mathfrak S_4$-torsor. Once
ordered, the unique coordinate taking the first three to
$\infty,0,1$ identifies the fourth with $\lambda\in U$; this
identification works in families, since a genus-zero curve with
three disjoint sections has that unique trivialization. Changing
the ordering gives the six cross-ratios
\[
 \lambda,\quad1-\lambda,\quad1/\lambda,\quad
 1/(1-\lambda),\quad(\lambda-1)/\lambda,\quad
 \lambda/(\lambda-1).
\]
Its kernel on the four labels is the double-transposition group.
Descent along the ordering torsor proves the stack equivalence,
including nontrivial conics over the base. In particular its generic
geometric object has four automorphisms, whereas an object represented
by a scheme has trivial automorphism group.

Set $t=j/256$ and $A=R[\lambda,\lambda^{-1},(1-\lambda)^{-1}]$.
The expression $t$ is fixed by both $\lambda\mapsto1-\lambda$ and
$\lambda\mapsto1/\lambda$, which generate the effective action.
The six displayed cross-ratios are the roots of the monic polynomial
\[
 (X^2-X+1)^3-tX^2(1-X)^2.
\]
This identity can be checked by substituting each of the six roots
in the fraction field, where they are generically distinct.
It follows that the three generators $\lambda,\lambda^{-1}$ and
$(1-\lambda)^{-1}$ are integral over $R[t]$, hence $A$ is finite over
$R[t]$. Over the fraction field $K=\mathbb Q(i)$, the rational map
$\lambda\mapsto t$ has degree~6, with coprime numerator and
denominator. The six automorphisms above therefore have fixed field
$K(t)$. Every element of $A^{\mathfrak S_3}$ lies in $K(t)$ and is
integral over $R[t]$. Since $R$ is a localization of the Gaussian
integers, $R[t]$ is integrally closed, so
$A^{\mathfrak S_3}=R[t]=R[j]$. This proves the stated quotient over
$R$. To pass from the invariant ring to the coarse moduli assertion
for algebraic-space targets, we apply Alper (2026), Theorem~5.2.3:
for any finite abstract group acting on an affine scheme,
$[\operatorname{Spec}A/G]\to\operatorname{Spec}A^G$ is a coarse
moduli space, with no assumption that $|G|$ is invertible.
The downloaded January~5 version and its proof on pp.~177--181
were opened. In a geometric fibre of odd characteristic the same degree-six
map is separable: in characteristic~3 its derivative is not zero,
and elsewhere this follows directly by differentiation. Thus it has
the required orbit fibres even at the wild prime~3. The groupoid
construction of Theorem~\ref{A:fourpoint-groupoid} is compatible with
these ordered charts and their descent data, giving its claimed
interpretation.
\end{proof}

\section{A.6. Anabelian assertions and mapping class presentations}

\paragraph{The group and the evidence must be specified.}
The original scans 67/048--051 were opened, including a rotated rendering
of 051. The handwritten date on 048 is only ``15.1.; the
reused printout is dated August~1982. The supplied study assigns
15~January~1984 from its surrounding dated material and vocabulary.
We retain that as a probable contextual date, not as a year written
on the sheet. The symbol $\pi_1([Y,N])$ in 048 is intended to retain the
topology of a smoothing; it cannot mean the ordinary geometric \'etale
fundamental group of the underlying proper nodal curve. We use
$\operatorname{Mod}_{g,n}$ for orientation-preserving mapping classes
permuting $n$ punctures, and $\operatorname{PMod}_{g,n}$ when every
puncture is fixed. A prime and an exclamation mark in the manuscript
must not be conflated: the upper-left of 051 has $\mathcal T_{0,4}^{+!}$,
whereas its upper-right has $\mathcal T_{0,4}^{+\prime}$. The supplied
transcription reads the latter as another $+!$. The two enlarged crops
in \texttt{qa/A567\_051\_topleft-1.png} and
\texttt{qa/A567\_051\_topright-1.png} resolve the difference on the
original. We give the correct groups below without guessing an
unwritten global convention for the prime.

\begin{proposition}[Refuted: the unqualified criterion in 67/049(b)]
\label{A:inertia-maximality}
Conjugacy classes of all closed subgroups isomorphic to
$\widehat{\mathbb Z}$, externally invariant under the arithmetic
Galois group and acted on cyclotomically, cannot be in bijection with
the finite set of nodes and marked points of a stable curve. This is
false even for the smoothing fundamental group of the split stable
four-pointed genus-zero curve consisting of two three-pointed lines.
For a smooth hyperbolic curve over a field with non-Tate
$\ell$-adic cyclotomic character, a correct cuspidal criterion uses
\emph{maximal} closed procyclic cyclotomically normalized subgroups
with nonzero $\ell$-primary part. It identifies the cusps, not every
possible nonperipheral cutting curve.
\end{proposition}

\begin{proof}
Take two copies of $\mathbb P^1_{\mathbb Q}$ with coordinates $z,w$,
identify $z=0$ with $w=0$, and mark $1,\infty$ on each component.
Both components have three special points, so the curve is stable.
Choose the rational normal tensor
$\partial_z\otimes\partial_w$ at its node. In its punctured
smoothing group, a loop about any one external marked point generates
a closed subgroup $C\simeq\widehat{\mathbb Z}(1)$. The arithmetic
Galois group preserves its conjugacy class and acts on it by the
cyclotomic character; the marked point and its local parameter are
rational. These properties also hold for
\[
 mC\subset C\qquad(m=1,2,3,\ldots).
\]
Multiplication by $m$ is an isomorphism from $\widehat{\mathbb Z}$
onto $m\widehat{\mathbb Z}$, so every displayed subgroup still has
the required abstract isomorphism type. The geometric smoothing is
a four-punctured sphere. Its abelianization is
$\widehat{\mathbb Z}^{3}$ and the chosen peripheral class is a
primitive vector $e_1$. The images of the subgroups are the distinct
submodules $m\widehat{\mathbb Z}e_1$. Conjugation acts trivially on
abelianization, so these give infinitely many different conjugacy
classes. They cannot parameterize five special points. The same
argument already works on the smooth curve
$\mathbb P^1-\{0,1,\infty\}$.

For the corrected assertion, Stix (2012), Definition~29,
Theorem~30 and Corollary~31, proves precisely that cuspidal inertia
groups are the maximal closed procyclic, cyclotomically normalized
subgroups with nontrivial $\ell$-part. His hypothesis ``non-Tate''
means that every Galois-equivariant map
$\mathbb Z_\ell(1)\to T_\ell A$ vanishes for every abelian variety
$A$ over the field. A cyclotomic normalizer is required to project
onto an open subgroup of the arithmetic Galois group. Thus finite
extensions, needed to rationalize the cusp, are allowed. Applying the
corollary to each hyperbolic normalization component identifies its
punctures, including the branches that will become nodes once their
pairing is separately reconstructed. This last qualification is
essential: the cusp theorem on a smooth curve does not supply that
pairing by itself.
\end{proof}

\paragraph{Arithmetic fields in the corrected cusp criterion.}
The non-Tate hypothesis holds over every field finitely generated over
$\mathbb Q$. Here is the usual reduction, to make the specialization
input explicit. Spread an abelian variety and its Tate module over a
nonempty smooth finite-type arithmetic base on which it has good
reduction and $\ell$ is invertible. A Galois-equivariant homomorphism
$\mathbb Q_\ell(1)\to V_\ell A$ is invariant under the fundamental
group of that base, since both representations are unramified there.
At any closed point with residue field of cardinality $q$ it would
intertwine multiplication by $q$ on the Tate line with Frobenius on
$V_\ell A$. The latter has eigenvalues of complex absolute value
$q^{1/2}$, whereas $q\ne q^{1/2}$. It must vanish. Tensoring an
integral homomorphism with $\mathbb Q_\ell$ is injective. This proves
the required non-Tate assertion from the Frobenius theorem for
abelian varieties; it is also the weights argument mentioned in
Stix's \S5.1. No analogous conclusion is being asserted over
$\overline{\mathbb Q}$ with its trivial Galois group.

\begin{proposition}[Proved: what the three topological types mean]
\label{A:three-types}
For an oriented surface $S_{g,n}$, the geometric kinds of nontrivial
simple closed curve are peripheral, nonseparating, and separating
nonperipheral. Some kinds may be absent. Within the separating kind,
the mapping class orbit is specified by the genera and puncture sets
on the two sides, up to exchanging the sides. Within the peripheral
kind, a pure mapping class orbit specifies the puncture. These kinds
are not three conjugacy classes in a fundamental group, nor three
conjugacy classes of homomorphisms
$\widehat{\mathbb Z}(1)\to\pi_1$.
\end{proposition}

\begin{proof}
Cutting a simple curve either keeps the surface connected, separates
off a once-punctured disc, or gives two other components. The first
case is nonseparating and the second is peripheral. In the third case
the two components have genera $g_1,g_2$ with $g_1+g_2=g$ and
puncture sets partitioning the original marks; each has an extra
boundary circle. The classification of oriented surfaces with marked
points gives a homeomorphism of the cut surfaces exactly when these
data agree. Boundary collars extend it across the seam. This proves
the stated orbit classification. For a labelled partition $(I_1,I_2)$
the nonperipheral condition on each side is
$2g_i+|I_i|\ge2$. A nonseparating curve exists only when $g\ge1$.
Thus already $S_{1,1}$ has only the peripheral and nonseparating
kinds, whereas many surfaces have several separating orbits.

For a direct contradiction to the literal ``three conjugacy
classes'' reading of the NB at 048, let
$\pi_1(S_{1,1})=F(a,b)$. For every integer $m\ge0$, the primitive
word $ab^m$ represents a simple nonseparating loop. It is primitive
because $(ab^m,b)$ is a free basis. Its closure in the profinite
completion is a copy of $\widehat{\mathbb Z}$: the projection onto
the first coordinate in the abelianization is a left inverse to
the homomorphism generated by $ab^m$. The corresponding
homomorphisms have abelianized values $(1,m)$, so no two are
conjugate. There are infinitely many. The word ``three'' is
therefore tenable only as a coarse classification into possible
topological kinds, with further types inside it, and not as the
literal cardinality of a set of conjugacy classes. The scan confirms
the words ``trois classes de conjugaison'', but does not define a
replacement equivalence relation that removes this distinction.
\end{proof}

\begin{proposition}[Proved obstruction: ordinary \'etale groups lose a node]
\label{A:ordinary-nodal-loss}
Let $Y$ be the proper split stable four-pointed genus-zero curve in
Proposition~\ref{A:inertia-maximality}. Its ordinary geometric
\'etale fundamental group is trivial. Removing its four external
marks gives a free profinite group of rank~2. A smooth four-punctured
genus-zero fibre has free profinite group of rank~3. Thus neither
ordinary group can serve as the smoothing group in 67/048.
\end{proposition}

\begin{proof}
Normalize $Y$ to two projective lines. A finite \'etale cover pulls
back to a disjoint union of trivial covers on each line, because a
nontrivial connected \'etale cover of $\mathbb P^1$ in characteristic
zero would violate Riemann--Hurwitz. Descent across the single node
identifies their two fibres by a bijection. Each matched pair is one
connected copy of $Y$, so a connected cover has degree one.

With the external marks removed, each normalized component is
$\mathbb G_m$ and has geometric fundamental group
$\widehat{\mathbb Z}$. The same description of a finite cover,
now by a finite set acted on by each of the two component groups
and an identification of their fibres at the node, is exactly the
category of finite sets with an action of their profinite free
product. It is the free profinite group on two generators. In
contrast, cutting the smooth four-punctured sphere along arcs to
three of its punctures gives the usual presentation with four
peripheral generators and their single product relation; it is
free on three. Their abelianizations have different ranks. This
also proves the failure without appealing to a nonabelian invariant.
\end{proof}

\paragraph{The proved logarithmic replacement for 049(a,b).}
We state exactly the available reconstruction input. A stable log
curve over the standard log point has local nodal charts of the form
$xy=t^e$, $e>0$, with $t$ the log parameter. Its geometric
logarithmic group is the fundamental group of a graph of groups:
vertices carry the groups of the punctured normalization components,
closed edges carry the cyclic node groups, and open edges carry
cuspidal groups. One retains the outer action of logarithmic base
inertia $I\simeq\widehat{\mathbb Z}^{\Sigma}(1)$, for a nonempty
set of primes $\Sigma$ invertible in the residue field.

\begin{theorem}[Known, with explicit derivation: logarithmic graphicity]
\label{A:log-graphicity}
For stable log curves over an algebraically closed standard log
point, an isomorphism of their geometric pro-$\Sigma$ groups that
intertwines their logarithmic inertia actions and preserves the
cuspidal inertia classes induces an isomorphism of the graphs of
groups. This includes recovery of the vertices, nodes, cusps, and
incidence relations. For curves without cusps the cuspidal condition
is empty. More generally, for the outer representations of
Hoshi--Mochizuki's Definition~2.4(iii), the same conclusion holds
under their NN hypothesis, nonempty cusp set, and group-theoretic
cuspidality.
\end{theorem}

\begin{proof}
For a geometric stable log curve, its base-inertia representation is
of IPSC type. Hoshi and Mochizuki (2011), Remark~2.4.2, proves the
implications IPSC $\Rightarrow$ SNN $\Rightarrow$ NN. When cusps
are present, apply their Corollary~4.2 with the given geometric groups
as $\Pi_{\mathcal G},\Pi_{\mathcal H}$ and the logarithmic inertia
groups as $I,J$. The intertwining condition is exactly its commuting
outer-action diagram; the other two hypotheses are NN and cuspidality.
Its conclusion ``graphic'' is preservation of the vertex and edge
subgroups with their incidence data. When no cusps are present,
Mochizuki (2007), Corollary~2.7(iii), applies directly: the full
logarithmic inertia image is logarithmically full, and the stipulated
isomorphism preserves that inertia. This establishes the first
assertion in all cases. The second assertion is precisely
Corollary~4.2 with its hypotheses retained; no geometric realization
of a general NN representation is inferred.
\end{proof}

\paragraph{The normal tensors and the arithmetic group in 049(a).}
The choice of group matters. We now give the tangential interpretation of
$\pi_1([Y,N])$, including the comparison that is needed to use the
reconstruction theorems. The standard log point over $K$ means
\[
 S_K^{\log}=(\operatorname{Spec}K,\ K^\times\oplus\mathbb N),
 \qquad t=(1,1),\quad t\longmapsto0.
\]
The element $t$, not just its characteristic monoid, is fixed. Its
logarithmic Galois group is
$G_K^{\log}=\widehat{\mathbb Z}(1)\rtimes G_K$. The specified splitting
$s_t:G_K\to G_K^{\log}$ fixes compatible formal roots $t^{1/n}$ and
acts in the usual way on $\overline K$. No logarithmic inertia group is
added to the data on the right side of Theorem~\ref{A:tangential-anabelian}
below.

\begin{lemma}[Proved comparison: tensors, log curves, and tangential groups]
\label{A:tensor-log-comparison}
Let $(Y,D)$ be a stable pointed curve over a characteristic-zero field $K$.
At a geometric node $e$, write $e^+,e^-$ for its two normalization
branches and $L_e=T_{e^+}\widetilde Y\otimes T_{e^-}\widetilde Y$.
A Galois-equivariant collection $N_e\in L_e\setminus\{0\}$ defines a
stable log curve $Y_N^{\log}\to S_K^{\log}$, with thickness one at every
node, and hence an outer action
\[
 \rho_N:G_K\longrightarrow\operatorname{Out}(\Delta_N),\qquad
 \Delta_N=\pi_1^{\log}(Y_N^{\log}/\overline S_K^{\log}).
\]
Here $\rho_N$ is the restriction along $s_t$ of the logarithmic base
action. Its geometric group is the profinite fundamental group of a
smooth surface of the same type as $(Y,D)$. Over $\mathbb C$ it is the
profinite completion of the topological surface obtained by gluing the
oriented real boundary circles at the nodes using the phases of $N_e$.
The construction depends only on $(Y,D,N)$, not on higher order
plumbing coordinates or a choice of smoothing family.

After a finite extension splitting the nodes, take two tensor systems
$N,N'$ on the same $Y$, and put $\delta_e=N'_e/N_e\in K^\times$.
Identify their geometric groups and graphs by choosing compatible roots
$\delta_e^{1/n}$ in $\overline K$. In the Dehn-twist coordinates of
Hoshi and Mochizuki (2012), Theorem~4.8, the difference of the two
outer actions is
\begin{equation}\label{A:node-Kummer-action}
 \rho_{N'}(g)=T^{\kappa(\delta)(g)}\rho_N(g),\qquad
 \kappa(\delta_e)(g)=
 \bigl(g(\delta_e^{1/n})/\delta_e^{1/n}\bigr)_n
 \in\widehat{\mathbb Z}(1).
\end{equation}
The orientation of each twist is chosen to agree with the Kummer
inertia of its smoothing coordinate.
\end{lemma}

\begin{proof}
We first explain precisely which normal line occurs. A split node has
completed ring $K[[x,y]]/(xy)$; its one-parameter smoothing has equation
$xy=s$. Replacing $x,y$ by $a x+\cdots,b y+\cdots$ changes $s$ to
$ab s+\cdots$. Thus the tangent vector in the smoothing direction
corresponds to $\partial_x\otimes\partial_y$. In particular
$N_e=q_e\partial_x\otimes\partial_y$ gives the logarithmic chart
$xy=q_e t$: in a chart $\mathbb N^2\to\mathcal O_Y$ given by $x,y$,
the base generator maps to the unit $q_e^{-1}$ times the sum of the
two generators. At a marked point use the usual divisorial log chart,
and away from marks and nodes pull back the base log structure.

The coordinate independence can also be checked without manipulating
formal branches. On an \'etale chart of the stable-curve moduli stack,
the boundary is a normal crossings divisor with one local equation
$s_e$ for each node. This is the moduli log structure used in
Hoshi and Mochizuki (2012), \S0, ``Curves''. Pullback to its $K$-point
gives the characteristic monoid $\mathbb N^{\{e\}}$. A log map from
$S_K^{\log}$ sending the generator for $e$ to $q_e t$ is exactly a
nonzero vector in the corresponding normal line. If
$s'_e=u_e s_e$, the coefficient changes to $u_e(Y)q_e$; this is the
normal-line transformation law. Units equal to one at $Y$ make no
change. The local identifications therefore satisfy the cocycle
condition on overlaps and descend, including when branches or nodes
are permuted. Pulling back the universal stable log curve gives the
asserted $Y_N^{\log}$. Equivalently, the familiar local deformation
calculation gives independent smoothing coordinates: the obstruction
group is zero because $H^2$ of the tangent sheaf of a curve and $H^1$
of its skyscraper local-deformation sheaf vanish. The node contribution
to the tangent space is precisely $\bigoplus_e L_e$ computed above.

For the passage to groups we use Mochizuki (2007), Example~2.5,
which constructs the admissible-cover exact sequence
\[
 1\longrightarrow\Delta_N\longrightarrow\Pi_{Y_N^{\log}}
 \longrightarrow G_K^{\log}\longrightarrow1
\]
and identifies $\Delta_N$ with the profinite fundamental group of the
associated semi-graph of anabelioids. Its vertices are the punctured
normalization components and its edges are their branch inertia,
identified inversely at a node. This is the usual van Kampen
presentation of the smoothing surface. Restriction along $s_t$ gives
$\rho_N$.

There is a direct comparison with the tangential construction, rather
than just an analogy of terminology. Hoshi, Matsumoto, and Nakayama
(2026), Proposition~1, identifies the tangential fibre functor with the
logarithmic root construction that adjoins all roots of the boundary
coordinates. Their Remark~1 proves its independence of coordinate
units equal to one on the boundary stratum; it depends only on the
normal cotangent basis, or equivalently the nonzero normal vectors.
Their Proposition~4 identifies its geometric groupoid with the
profinite completion of the analytic groupoid. These are assertions
of the opened June~2026 preprint, not a claim of a later journal
publication. They apply to each \'etale scheme chart above; functoriality
and descent give the same statement on the moduli stack. In the local
analytic chart $xy=q_e t$, the real oriented boundary equation is
$\arg x+\arg y=\arg q_e+\arg t$. On the fibre $\arg t=0$ this is
exactly the anti-oriented circle gluing specified by $N_e$. Its modulus
does not change the topological group, but it does enter the arithmetic
tangential functor.

It remains to compute that arithmetic dependence. Work first at one
node in a split chart. In the Kummer covering of the smoothing
coordinate, choose $s^{1/n}$. The tangential section with coefficient
$q$ fixes $(s/q)^{1/n}$, so it sends $s^{1/n}$ to
$g(q^{1/n})q^{-1/n}s^{1/n}$. Thus replacing $q$ by $q'$ changes the
section by the Kummer cocycle of $q'/q$. Equivalently the local root
equation is
\[
 x^{1/n}y^{1/n}=q^{1/n}t^{1/n};
\]
fixing $t^{1/n}$ leaves exactly the displayed discrepancy in the two
branch identifications. For several nodes the boundary chart is a
product of these Kummer coordinates, so the discrepancies are
independent and form the vector $(\kappa(\delta_e))_e$.

Finally this local base inertia is the asserted geometric multitwist,
with no unnoticed kernel or rescaling: Hoshi and Mochizuki (2012),
Theorem~5.7, identifies the universal boundary inertia with
$\bigoplus_e\Lambda_{\mathcal G}$ under the isomorphism
$\operatorname{Dehn}(\mathcal G)\simeq\bigoplus_e\Lambda_{\mathcal G}$
of their Theorem~4.8(iv). It also identifies the scheme-theoretic
Kummer cyclotomes with these combinatorial cyclotomes. Its hypotheses
are checked after base change to $\overline K[[t]]$: this is a complete
discrete valuation ring with algebraically closed characteristic-zero
residue field, so every prime is allowed. A smoothing with all
coefficients $q_e$ exists by the unobstructed local deformation
calculation above. The comparison is functorial under its Galois
symmetries and therefore descends. This proves
\eqref{A:node-Kummer-action}. Changing any chosen roots changes the
cocycle by the corresponding coboundary and the identification by the
same multitwist, as required.
\end{proof}

\begin{lemma}[Proved: full Kummer faithfulness in the required fields]
\label{A:Kummer-faithfulness}
If $K$ is finitely generated over $\mathbb Q$, the Kummer map
\[
 K^\times\longrightarrow H^1(G_K,\widehat{\mathbb Z}(1))
\]
is injective, and $\widehat{\mathbb Z}(1)^{G_K}=0$.
\end{lemma}

\begin{proof}
A zero class has zero image in every $H^1(G_K,\mu_n)$, so the ordinary
Kummer exact sequence shows that its representative $a$ is an $n$th
power in $K$ for every $n$. If $a$ were transcendental over $\mathbb Q$,
include $a$ in a transcendence basis of $K/\mathbb Q$. The discrete
valuation at $a=0$ of the resulting rational function field extends to
a discrete valuation $v$ of its finite extension $K$, with
$0<v(a)<\infty$. Being an $n$th power for every $n$ would make every
$n$ divide $v(a)$, a contradiction. Thus $a$ belongs to the number
field $F$ of elements of $K$ algebraic over $\mathbb Q$. Every $n$th
root of $a$ lying in $K$ is also in $F$.

Valuations at primes of $F$ first show that $a$ and all these roots
are units. Their conjugate absolute values are uniformly bounded,
since if $b^n=a$ then $|\tau(b)|=|\tau(a)|^{1/n}$. There are only
finitely many algebraic integers in $F$ with this bound: their monic
characteristic polynomials over $\mathbb Q$ have bounded integral
coefficients, so only finitely many polynomials and roots can occur.
An element $b$ must consequently repeat for two distinct exponents;
then $b^{n-m}=1$, hence $a$ is a root of unity. The group of roots of
unity in $F$ is finite by the same bounded-coefficient argument.
Taking $n$ divisible by its order forces $a=1$. This proves injectivity
without a finite-generation assertion for $K^\times$.

For the invariants, for each prime $\ell$ the image of
$\chi_\ell:G_K\to\mathbb Z_\ell^\times$ is open: the constants of
$K$ form the finite extension $F/\mathbb Q$, and regularity over $F$
makes $G_K\to G_F$ surjective. Thus some element has
$\chi_\ell(g)\ne1$. If $u\in\mathbb Z_\ell(1)$ is invariant,
$(\chi_\ell(g)-1)u=0$, and torsion-freeness implies $u=0$.
This holds in every $\ell$-component of $\widehat{\mathbb Z}(1)$.
\end{proof}

\begin{theorem}[Proved from the cited reconstruction inputs: 049(a)]
\label{A:tangential-anabelian}
Let $K$ be finitely generated over $\mathbb Q$. Let
$(Y,D,N)$ and $(Y',D',N')$ be geometrically connected stable pointed
curves of hyperbolic type with a nonzero tensor at every node, and
use the tangential groups of Lemma~\ref{A:tensor-log-comparison}.
The natural map
\[
 \operatorname{Isom}_K((Y,D,N),(Y',D',N'))
 \longrightarrow
 \operatorname{Isom}_{G_K}^{\mathrm{out}}
 ((\Delta_N,\rho_N),(\Delta_{N'},\rho_{N'}))
\]
is a bijection. On the right, isomorphisms are continuous and considered
modulo geometric inner automorphisms. If markings are labelled,
require their peripheral conjugacy classes to match their labels.
In particular the arithmetic outer action recovers the curve, its
node and cusp incidence, and the actual nonzero tensors, rather than
only their orders or their classes modulo powers.
\end{theorem}

\begin{proof}
\emph{Recovery of the graph.}
We check the hypothesis of Mochizuki (2007), Corollary~2.7(ii),
instead of imposing preservation of logarithmic inertia that is not
present in the right side. For any prime $\ell$, the action of $G_K$
on the semi-graph $\mathcal G$ of Lemma~\ref{A:tensor-log-comparison}
is $\ell$-graphically full in the precise sense of his
Definition~2.3(iii). The cyclotomic image is open by
Lemma~\ref{A:Kummer-faithfulness}. For each characteristic finite
cover of $\mathcal G$, the module
\[
 (M^{\mathrm{vert}}/M^{\mathrm{edge}})\otimes\mathbb Z_\ell
\]
is the sum of the Tate modules of the Jacobians of the proper
normalization components; this is the geometric description in
Mochizuki's Proposition~1.3 and its proof. After passage to an open
subgroup the finite deck and component permutations disappear, and
these finitely many Jacobians and covers are defined over a finite
extension $K'/K$.
More explicitly, a characteristic geometric cover need not descend
with the original $G_K$-action; choose its arithmetic descent after a
finite extension and then restrict to the finite-index subgroup
$\Pi_{\mathrm{cover},K'}\subset\Pi_{\mathcal G}^{G_K}$.
Its action on the cover's homology factors through $G_{K'}$, since
the geometric kernel acts by inner automorphisms; associated weights
and quasi-trivial ranks may be computed after precisely this
finite-index restriction, so the finite deck ambiguity causes no
additional weights.

To check the associated weights, spread these data over a nonempty
normal scheme of finite type over $\mathbb Z[1/\ell]$, shrinking so
that the Jacobians have good reduction. At a closed point with finite
residue field of cardinality $Q$, their Frobenius eigenvalues have
complex absolute value $Q^{1/2}$. This is the same weight-one input
used in the proof of Mochizuki's Proposition~1.3. If a rational
cyclotomic character $\psi$ contributes an associated weight,
$\psi^b=\chi_\ell^a$ after a finite extension, its Frobenius absolute
value is $Q^{a/b}$. Consequently $2a/b=1$. The argument works for
every cover and every open subgroup used in the definition; hence
all associated weights of the displayed quotient belong to
$\{1\}\subset(0,2)$. The required fullness follows. Notice that
there is just the one original nodal degeneration: the spreading was
chosen to preserve the smooth normalization components and their
Jacobian good reduction.

Apply Corollary~2.7(ii) to the given outer isomorphism $\alpha$, with
$J_{\mathcal G}=J_{\mathcal H}=G_K$ and the identity isomorphism of
these groups. It is graphic. Thus it identifies the vertices, their
punctured-curve groups, the cusps, the nodes, and the two branches of
each node. This argument also handles a smooth curve, whose graph
has one vertex and no nodes.

\emph{Recovery of the normalization and its gluing.}
Pass temporarily to a finite extension splitting the vertices, nodes,
and branches. The restricted vertex outer actions are exactly the
usual arithmetic outer actions on the smooth normalization components
with all marks and node branches deleted; they are independent of
$N$. This follows directly either by restricting the log charts away
from nodes or by restricting the admissible-cover description in
Mochizuki's Example~2.5. Each component is hyperbolic by stability.
Its geometric fundamental group is centre-free. Consequently an
isomorphism of its outer $G_{K'}$-actions gives the corresponding
isomorphism of arithmetic extensions: an extension with centre-free
kernel is canonically the pullback of
$\operatorname{Aut}(\Delta)\to\operatorname{Out}(\Delta)$
along the prescribed outer action.
The field is sub-$p$-adic: it embeds in a finite extension of a
rational function field over $\mathbb Q_p$ for any chosen prime $p$.
Mochizuki (1999), Theorem~A, therefore applies to its full geometric
fundamental group. Apply the theorem to each induced open group
isomorphism and to its inverse. The two resulting dominant maps are
inverse because the theorem is injective on maps. This gives unique
isomorphisms of the punctured normalization components. They extend
uniquely to their smooth proper compactifications and match the
punctures dictated by graphicity. Matching the paired branches glues
them to an isomorphism $f:(Y,D)\to(Y',D')$. Uniqueness on the components
makes $f$ Galois equivariant and gives descent to $K$.

\emph{Recovery of the tensors.}
Transport $N'$ back by $f$, so that $Y'=Y$ and the induced
isomorphisms on all vertices are the identity. Use the geometric
identification and root choices of
Lemma~\ref{A:tensor-log-comparison}. By the definition of the
Dehn-twist kernel, $\alpha$ now lies in
$\operatorname{Dehn}(\mathcal G)$. Hoshi and Mochizuki (2012),
Theorem~4.8(iv), gives
\[
 \operatorname{Dehn}(\mathcal G)
   \simeq\bigoplus_{e\in\operatorname{Node}(\mathcal G)}
                 \widehat{\mathbb Z}(1).
\]
Write $\alpha=T^u$. Their Theorem~4.8(v) says that Galois conjugation
on this module permutes the nodes and multiplies by the cyclotomic
character. Over the splitting field there is no permutation. By
\eqref{A:node-Kummer-action}, the equivariance equation is therefore
\[
 \kappa(\delta_e)(g)=(1-\chi(g))u_e\qquad(g\in G_{K'}).
\]
Thus the Kummer class of every $\delta_e$ is zero. Since $K'$ is still
finitely generated over $\mathbb Q$, Lemma~\ref{A:Kummer-faithfulness}
gives $\delta_e=1$. Hence $f$ preserves every actual tensor.
The assertion descends from $K'$ because it is an equality of
geometric tensors and $f$ is already defined over $K$.

\emph{Uniqueness of the group isomorphism.}
Once $f$ preserves $N$, any difference between $\alpha$ and its induced
outer isomorphism is a Galois-invariant multitwist. After the same
finite splitting extension, each of its coordinates belongs to
$\widehat{\mathbb Z}(1)^{G_{K'}}=0$. The difference is zero.
Conversely any isomorphism of the decorated curves induces a compatible
outer isomorphism by functoriality of the tangential construction.
This proves both surjectivity and injectivity of the stated map.
\end{proof}

\begin{example}[Proved: the smallest tensor test, completely determined]
\label{A:nodal-q-test}
Take two labelled projective lines, each with node at $0$ and external
marks at $1,\infty$, over $\mathbb Q$. Set
$N=q\,\partial_z\otimes\partial_w$, $q\in\mathbb Q^\times$.
The labelled underlying curve has no automorphisms. Its geometric
tangential group is $\widehat F_3$, presented as the amalgam of the
two three-punctured-sphere groups along their inverse node generators.
Writing it as $\widehat F(c,a,b)$, a twist fixes the left subgroup
$\langle c,a\rangle$ and conjugates the right subgroup
$\langle c,b\rangle$ by $c$. The outer actions for $q,q'$ are
isomorphic over $G_{\mathbb Q}$ if and only if $q=q'$. Indeed any
such isomorphism is graphic and identity on the labelled components,
so it is $T^u$; the necessary and sufficient equation is
$\kappa(q'/q)=(1-\chi)u$, which forces $q'/q=1$.
Thus the bare geometric group loses $q$, but its specified arithmetic
tangential action does not. This rules out using the former
independence as a counterexample to 049(a).
\end{example}

\paragraph{The exact extent of the answer to 049(b).}
Theorem~\ref{A:tangential-anabelian} in particular recovers the finite
set of node and cusp conjugacy classes, with incidence, from the
arithmetic outer action. This proves the intended reconstruction
claim for the tangential group. It does not make the literal
criterion using \emph{all} cyclotomic procyclic subgroups true:
Proposition~\ref{A:inertia-maximality} remains a counterexample.
For smooth curves, the maximality version quoted there is Stix's
precise theorem. In the nodal case the fully justified replacement
here is the graphic reconstruction just proved; no unproved claim
that merely adding the word ``maximal'' is a sufficient singular
criterion is being made.

\paragraph{Sources and corrections for the reconstruction.}
The derivation uses Mochizuki (2007), Definition~2.3,
Corollary~2.7(ii), and Example~2.5; Mochizuki (1999), Theorem~A;
Hoshi and Mochizuki (2012), Definition~4.4, Proposition~4.5,
Theorem~4.8(iv--v), pp.~78--82, and Theorem~5.7, pp.~105--106
of the author version; and the explicit
tangential comparison of Hoshi, Matsumoto, and Nakayama (2026),
Propositions~1 and 4 and Remark~1. The author's 2019 and 2020
comments on the first two papers and all of the January~2024
comments on the 2012 paper were opened. In particular the corrected
meaning of $\operatorname{Aut}^{|\mathrm{grph}|}(\mathcal G)$ is
identity on the underlying semi-graph, exactly as used in the
Dehn-twist step. None of these comments removes the stated
reconstruction conclusions. Mochizuki (1996), \S6, Propositions~6.1
and 6.4, and Mochizuki (2007), Remark~2.8.1, provide the earlier
logarithmic decomposition-group route; the proof above supplies the
extra tangential-section and Kummer steps explicitly rather than
identifying the two routes without proof.

\paragraph{049(c) and the limit of Lead L6.}
The scan asks for a definition of an anabelian compactification of an
arbitrary smooth anabelian variety, inspired by Mumford--Deligne
compactification, and asks how it behaves under specialization.
The curve theorem above gives a precise model using stable log curves.
There is also a proved higher-dimensional model for saturated
hyperbolic polycurves, given below; its hypotheses cannot be replaced
by the undefined condition ``arbitrary anabelian variety''. The
distinction already matters for a product of two three-punctured lines.

\begin{proposition}[Proved obstruction for 049(c): a compactification is extra data]
\label{A:compactification-obstruction}
Even for a known smooth anabelian surface over $\mathbb Q$, its
arithmetic fundamental group does not determine an arbitrary chosen
smooth projective compactification with simple normal crossings
boundary. In particular such compactifications cannot form a category
on which the functor assigning only the fundamental group of the
interior is fully faithful on isomorphisms.
\end{proposition}

\begin{proof}
Put $C=\mathbb P^1_{\mathbb Q}\setminus\{0,1,\infty\}$ and
$U=C\times C$. This is a strongly hyperbolic Artin neighbourhood:
the tower $C\times C\to C\to\operatorname{Spec}\mathbb Q$ has
elementary hyperbolic-curve fibrations, and every stage embeds in a
product of hyperbolic curves. Thus Schmidt and Stix (2016),
Definition~5.7 and Corollary~1.6 in the opened author version, apply
and give its anabelian isomorphism theorem in terms of the full
arithmetic fundamental group. This verifies the adjective
``anabelian'' for the example rather than assuming it.

One compactification is
$\overline U=\mathbb P^1\times\mathbb P^1$, with its six horizontal
and vertical boundary curves. Blow up their rational intersection
$(0,0)$ and put
$\overline U'=\operatorname{Bl}_{(0,0)}\overline U$, with the reduced
total transform $D'$ of the boundary. The blowup is smooth and
projective, and is the identity over $U$. Its boundary is again simple
normal crossings: in the two local charts $y=xu$ and $x=yv$ the
reduced total transform of $xy=0$ is respectively $xu=0$ and $yv=0$.
The pair $(\overline U',D')$ has seven irreducible boundary components,
including the exceptional curve, whereas $(\overline U,D)$ has six.
They are not isomorphic as pairs, even geometrically. Nevertheless
their interiors are the identical $\mathbb Q$-variety $U$, so their
arithmetic interior fundamental groups, including the augmentation
to $G_{\mathbb Q}$, are identical. This proves both assertions.
The finite component lists are also recorded by
\texttt{code/A567\_matrices.py}; the local blowup calculation, not
that enumeration, proves the geometric assertion.
\end{proof}

\begin{theorem}[Known inputs, complete deduction: canonical polystable compactifications]
\label{A:polystable-compactification}
Let $K$ be finitely generated over $\mathbb Q$. Consider strongly
hyperbolic Artin neighbourhoods $U/K$ that admit a hyperbolic-curve
tower whose stages are split and saturated in the sense of
Mochizuki (1999, \emph{Extending families}), Definitions~7.2--7.3.
In characteristic zero it suffices that the punctures of each stage
are disjoint sections and its proper smooth curve admits a full
level-$\ell$ structure for some prime $\ell\ge3$.
Then $U$ has a canonical proper log-polystable compactification
$X_U^{\log}$, independent of the choice of a tower. Moreover
\[
 \operatorname{Isom}_K(X_U^{\log},X_V^{\log})
 \xrightarrow{\ \sim\ }
 \operatorname{Isom}^{\mathrm{out}}_{G_K}
       (\pi_1(U),\pi_1(V)).
\]
Here the arithmetic groups retain their augmentations to $G_K$,
and the outer quotient is by the geometric inner automorphisms.
The statement includes $U=(\mathbb P^1_K\setminus\{0,1,\infty\})^d$
for every $d\ge1$.
\end{theorem}

\begin{proof}
The compactification input is Mochizuki's Corollary~7.4, author
pp.~29--31: a saturated polyhyperbolic interior over a log-regular
base has a log-polystable compactification. Corollary~6.2, author
pp.~27--28, extends an interior map from a log-regular pair to a
polystable pair uniquely. These exact statements, their hypotheses,
and their proofs were opened. Applying the extension result to
the identity between two compactifications and to its inverse
makes their extensions inverse, by uniqueness. This proves
canonical independence, including independence of a chosen tower.

Schmidt--Stix (2016), Corollary~1.6, reconstructs each interior
isomorphism uniquely from the indicated arithmetic group
isomorphism. Extend this interior isomorphism and its inverse by
the preceding result. Their extensions are inverse and are log
morphisms: the log structure is the divisorial one
$\mathcal O_X\cap j_*\mathcal O_U^\times$, which an
interior-preserving isomorphism preserves. Conversely, every log
isomorphism preserves the locus where the log structure is trivial,
so restricts to an interior isomorphism. This proves the bijection.

For the displayed products, use the coordinate-projection tower.
Every stage is the trivial three-punctured projective line; its
three punctures are sections and its Jacobian is zero, so the
level condition is automatic. Each stage embeds in a product of
hyperbolic curves, which verifies the strongly hyperbolic Artin
condition. Its compactification is $(\mathbb P^1)^d$ with the
coordinate boundary divisors. This pair is visibly polystable,
so uniqueness identifies it with the one just constructed.
\end{proof}

\paragraph{Definition supplied for the proved scope of 049(c).}
Within this class, an \emph{anabelian polystable compactification}
means the proper log-polystable pair whose trivial-log locus is
the given interior; its morphisms are the log morphisms preserving
that locus. ``Polystable'' requires a tower of proper stable
pointed curves, smooth over each preceding interior, with boundary
the union of the pulled-back marked sections and the base boundary
(Mochizuki, Definition~4.5). This condition excludes arbitrary
boundary blowups. No tower or level basis is retained as extra
data in the isomorphism theorem. Section~7, Remark~1 of the same
source shows that characteristic-zero hyperbolic polycurve towers
acquire saturation after suitable finite \emph{\'etale covers};
that is a statement about covers, not an assertion that the
original interior has been left unchanged.

\begin{theorem}[Proved from opened inputs: canonical normal compactification]
\label{A:polycurve-coarse-compactification}
Let $K$ be a field of characteristic zero and let $U/K$ be a
geometrically connected hyperbolic polycurve: $U$ admits a finite
tower of hyperbolic curves, in the sense of Mochizuki (1999),
Definition~4.6. No splitness or level structure is assumed.
There is a canonical normal projective pair $(\overline U,U)$
over $K$, obtained by the following procedure:
\begin{enumerate}
\item Choose a connected finite \'etale cover $V\to U$ which,
over a finite extension $K'/K$, has a split saturated hyperbolic
curve tower. Let $(X_V,V)$ be its polystable compactification.
\item Choose a connected finite \'etale Galois cover $W\to U$
dominating $V$, and normalize $X_V$ in $K(W)$, obtaining $Z$.
\item Extend $G=\operatorname{Gal}(W/U)$ to $Z$ and set
$\overline U=Z/G$.
\end{enumerate}
The resulting pair is independent, up to the unique isomorphism
which is the identity on $U$, of all covers, towers, and field
extensions in the procedure. Every $K$-isomorphism $U\simeq U'$
extends uniquely to an isomorphism of these pairs. If $U$ already
has a polystable compactification, this procedure returns it.

For $K$ finitely generated over $\mathbb Q$ and $U,U'$ strongly
hyperbolic Artin neighbourhoods, consequently
\[
 \operatorname{Isom}_K\bigl((\overline U,U),(\overline {U'},U')\bigr)
 \simeq
 \operatorname{Isom}^{\mathrm{out}}_{G_K}
       \bigl(\pi_1(U),\pi_1(U')\bigr).
\]
The right side retains the augmentations to $G_K$, with geometric
inner automorphisms divided out. The normal coarse boundary is not
asserted to be log regular; the log-regular normalization $Z$ is
retained whenever a log specialization theorem is applied.
\end{theorem}

\begin{proof}
We give the descent and independence argument because the
saturation hypothesis cannot simply be deleted from the
polystable compactification theorem.

\emph{Existence of the covering data.}
Mochizuki (1999, \emph{Extending families}), \S7, Remark~1
\cite{A567_Mochizuki1999Extending},
author p.~31, states that in characteristic zero finite \'etale
covers of hyperbolic polycurves remain hyperbolic polycurves,
and that a finite \'etale cover of the source and base makes
the tower split and saturated. Apply it with base $\operatorname{Spec}K$
and choose a connected component. Its new base is
$\operatorname{Spec}K'$ for a finite separable extension.
Corollary~7.4, author pp.~30--31, then produces $(X_V,V)$.
The successive stable-curve morphisms are projective because
their pointed relative dualizing sheaves are relatively ample
(Definition~4.3), so $X_V$ is projective over $K'$, and hence
over $K$. Definition~4.5 and Lemma~4.2 show that $(X_V,V)$
is a toric pair, that is, its divisorial log structure is log regular.
A connected finite \'etale cover of the connected normal scheme
$U$ is dominated by a connected finite \'etale Galois cover:
take the normal core of its open subgroup in the fundamental
group, or equivalently its Galois closure in the finite \'etale
cover category. This gives $W$.

\emph{Normalization and extension of the deck transformations.}
The normalization $Z\to X_V$ is finite, since $X_V$ is of
finite type over a field; thus $Z$ is normal and projective.
It restricts to $W\to V$. Every boundary valuation has residue
characteristic zero, so this finite cover is tamely ramified
there. Mochizuki's log purity theorem, Theorem~3.3, author
pp.~16--17, applies and says that $(Z,W)$ is log regular and
that its map to $(X_V,V)$ is log \'etale.

For $g\in G$ consider the map $W\xrightarrow{g}W\to V$.
It extends to a map $h_g:Z\to X_V$ by Corollary~6.2,
author pp.~27--28. Here one small base-field issue matters:
$g$ need not fix $K'$. Give $Z$ the $K'$-structure defined
by the embedding $K'\to K(W)$ induced by this map on $W$.
It extends across $Z$, since each of these algebraic constants
is integral over $K$ and hence lies in every normal local ring
of $Z$; the inverse of a nonzero constant has the same property.
Thus Corollary~6.2 is legitimately applied over $K'$.

The map $h_g$ lifts uniquely to an endomorphism $\widetilde g$
of $Z$. Indeed $Z$ is the normalization of $X_V$ in $K(W)$.
On any affine open of $X_V$, an element of its integral closure
in $K(W)$, transported by $g$, is integral over the pulled-back
coordinate ring. It is therefore regular in every local ring
of the normal source $Z$. These local lifts agree on $W$ and
glue. Applying the same argument to $g^{-1}$ gives an inverse:
the two compositions are the identity on the dense open $W$,
hence on $Z$, since $Z$ is separated. The same dense-open
argument proves $\widetilde g\widetilde h=\widetilde{gh}$.
Thus the actual group $G$, not only its individual elements,
acts on $Z$ and preserves $W$.

\emph{The projective normal quotient.}
For completeness, the finite-group quotient used here is a
scheme and is projective. If $A$ is an ample invertible sheaf
on $Z$, then $L=\bigotimes_{g\in G}g^*A$ is ample and has its
natural $G$-linearization, obtained by permuting the factors.
An ample section can be chosen not to vanish on a given finite
orbit: use a sufficiently high power to separate the finitely
many closed points, and use that $K$ is infinite. The product
of its $G$-translates is an invariant section. These invariant
sections give a cover of $Z$ by invariant affine opens.
On such an open $\operatorname{Spec}B$ take
$\operatorname{Spec}B^G$. The polynomial
$\prod_{g\in G}(T-gb)$ makes every $b\in B$ integral over
$B^G$; finite generation over $K$ then makes $B$ finite over
$B^G$. To also verify finite generation of the invariant ring,
let $B_0$ be the $K$-algebra generated by the coefficients of
these monic polynomials for a finite algebra generating set of
$B$. Then $B$ is a finite $B_0$-module, $B_0$ is noetherian,
and its submodule $B^G$ is finite over $B_0$. Invariant
principal localization commutes with taking
invariants, so these affine quotients glue. They identify
geometric points precisely by their finite $G$-orbits and
are categorical for invariant maps. Equivalently the glued
scheme is the Proj of the invariant section ring of $L$:
the invariant affine opens just constructed are its standard
charts. That ring is finitely generated (the section ring is
finite over its invariant subring), so the quotient is projective.
Each $B^G$ is normal because $B$ is normal: an element of
$\operatorname{Frac}(B^G)$ integral over $B^G$ is integral over
$B$, lies in $B$, and is invariant. Thus $\overline U=Z/G$
is normal. The invariant open $W$ descends to its finite
\'etale quotient $W/G=U$; its complement is closed because
the quotient is finite. This proves that $U$ is an open dense
subscheme of the projective quotient, with its original
$K$-structure.

\emph{Independence with one common covering field.}
Suppose first that $V_1,V_2$ are two saturated covers and one
connected finite \'etale Galois cover $W\to U$ dominates both.
Let $Z_i$ be the normalization of $X_{V_i}$ in $K(W)$.
Both pairs $(Z_i,W)$ are log regular by the same log purity
argument. The interior map $W\to V_2$ extends to
$Z_1\to X_{V_2}$ by Corollary~6.2, including the constant-field
argument above. Normality then lifts it to $Z_1\to Z_2$.
Reversing the roles gives $Z_2\to Z_1$. These maps restrict
to the identity on $W$, so are inverse. The resulting
isomorphism is unique and $\operatorname{Gal}(W/U)$-equivariant
by dense-open equality. Their quotients therefore coincide.
This also proves independence of the chosen towers.

\emph{Independence when the Galois cover changes.}
Any two finite \'etale Galois covers used above have a common
connected finite \'etale Galois cover $W'\to U$. It remains
to compare $W'$ with one of them, $W$. Write
$H=\operatorname{Gal}(W'/W)$ and normalize the same $X_V$
in these two function fields, obtaining $Z'$ and $Z$.
The map $Z'\to Z$ is finite. On each affine of $X_V$, the
$H$-invariants of the integral closure in $K(W')$ are
exactly the integral closure in $K(W)$: invariance imposes
the latter fraction field, and integrality over the same
base ring is unchanged. Thus $Z'/H=Z$. Taking the remaining
quotient proves
\[
 Z'/\operatorname{Gal}(W'/U)
   = Z/\operatorname{Gal}(W/U).
\]
Together with the preceding paragraph this proves independence
of all choices, including $K'$.

\emph{Canonicity and the arithmetic conclusion.}
Transporting a set of covering data along an isomorphism
$U\simeq U'$ and applying the proved independence yields its
extension. Uniqueness follows from equality on a dense open.
If a polystable $(X,U)$ already exists, normalize $X$ in
$K(W)$. Log purity makes that normalization log regular.
The common-field comparison with $X_V$ identifies it with
$Z$. Its $G$-quotient is $X$, by the same invariant
normalization argument with fixed field $K(U)$ and the
normality of $X$. Therefore the construction returns $X$.
Finally Schmidt--Stix (2016), Corollary~1.6, recovers the
interior isomorphism from exactly the indicated arithmetic
fundamental-group isomorphism for strongly hyperbolic Artin
neighbourhoods. Extend it by the canonicity just proved;
restriction is the inverse operation. This gives the stated
bijection of isomorphism sets.\qedhere
\end{proof}

\begin{definition}[Proved well-defined: an intrinsic compactification category for 049(c)]
\label{A:anabelian-compactification-definition}
Let $K$ have characteristic zero and let $U/K$ be a smooth
geometrically connected variety. A \emph{compactification with
the toric extension property} is a normal proper $K$-scheme $Y$
with an identified dense open $U\subset Y$, satisfying the
following two conditions.
\begin{enumerate}
\item There is a finite dominant morphism $Z\to Y$, where $Z$
is normal and integral, $K(Z)/K(Y)$ is Galois with finite group $G$, the
inverse image $W$ of $U$ is a finite \'etale $G$-torsor over
$U$, and $(Z,W)$ is a toric pair. Equivalently, the divisorial
fine saturated log structure on $Z$ has trivial locus $W$ and
is log regular. The existence of this cover is a property;
no choice of it is part of the compactification data.
\item For every excellent noetherian toric pair $(T,T^\circ)$
over $K$, every $K$-morphism $T^\circ\to U$ extends to a
$K$-morphism $T\to Y$.
\end{enumerate}
When the interiors belong to an anabelian isomorphism category
over $K$, meaning that their $K$-isomorphisms are recovered
bijectively from their arithmetic fundamental groups with the
$G_K$-augmentations, we call these its \emph{anabelian compactifications}.
The interior and its embedding are retained; morphisms are
morphisms of the pairs. Neither smoothness of $Y$ nor a log
regular structure on its coarse boundary is imposed.
\end{definition}

\begin{theorem}[Proved: intrinsic uniqueness and the toric extension property]
\label{A:intrinsic-anabelian-compactification}
The normal compactification of
Theorem~\ref{A:polycurve-coarse-compactification} satisfies
Definition~\ref{A:anabelian-compactification-definition}.
For an arbitrary smooth $U/K$, a compactification satisfying
that definition is unique if it exists. Every $K$-morphism
between two interiors having such compactifications extends
uniquely to their compactifications. In particular arithmetic
isomorphism reconstruction of an interior implies arithmetic
isomorphism reconstruction of this pair, without retaining
any auxiliary cover.
\end{theorem}

\begin{proof}
We first spell out a finite-morphism extension fact, including
the case of nondominant maps. Let $P\to X$ be finite, $S$ a
normal noetherian scheme, and $S^\circ\subset S$ an open set
dense in every component. Suppose given $S\to X$ and a lift
$S^\circ\to P$ over $X$. Then this lift extends uniquely to
$S\to P$. To prove it, work over an affine
$X_0=\operatorname{Spec}A$ with inverse image
$P_0=\operatorname{Spec}B$, where $B$ is finite over $A$,
and an integral normal affine $S_0=\operatorname{Spec}R$
mapping into $X_0$. A normal scheme has disjoint irreducible
components, so such affines cover $S$. Every $b\in B$ pulls
back on the dense open $S_0\cap S^\circ$ to a rational
function in $\operatorname{Frac}R$. Its monic integral
equation over $A$ becomes a monic equation over $R$. Since
$R$ is integrally closed, this rational function belongs to
$R$. The resulting homomorphism $B\to R$ extends the given
one, and the local homomorphisms agree on overlaps by equality
on a dense open. This argument does not assume that any map
to $X$ or $P$ is dominant.

Now fix the covering data $Z\to Y=\overline U$ and
$W\to U$ from Theorem~\ref{A:polycurve-coarse-compactification}.
The first condition of the definition holds by construction.
Given an excellent noetherian toric pair $(T,T^\circ)$ and a
map $a:T^\circ\to U$, form the finite \'etale $G$-torsor
\[
 T'^\circ=T^\circ\times_U W\longrightarrow T^\circ.
\]
Normalize each component of $T$ in this finite \'etale
algebra to obtain $T'\to T$. Excellence makes the
normalization finite. The algebra may be a product of fields;
normalization and the following arguments are applied to all
components. The $G$-action extends uniquely to $T'$.
On an affine normal component of $T$, the invariant subalgebra
of this normalization is the original coordinate ring:
its invariant generic algebra is the original function field,
and integrality plus normality gives the assertion. Thus
$T'/G=T$ as a finite categorical quotient.

Because $K$ has characteristic zero, ramification along each
boundary valuation is tame. The opened log purity theorem,
Mochizuki (1999), Theorem~3.3, makes $(T',T'^\circ)$ a
toric pair. The interior map $T'^\circ\to W\to V$ extends
to $b:T'\to X_V$ by Corollary~6.2. If $V$ is defined over
$K'$, the map to $\operatorname{Spec}K'$ first extends to
$T'$: the algebraic constants are integral over $K$ in each
normal local ring, exactly as in the preceding proof.
The finite-morphism extension fact just proved, applied to
$Z\to X_V$ and the given lift $T'^\circ\to W\subset Z$,
then supplies $\widetilde b:T'\to Z$. This step also handles
constant maps and all other nondominant maps.

For every $g\in G$, the maps
$\widetilde b\circ g$ and $g\circ\widetilde b$ agree on
$T'^\circ$. The source is reduced and the target separated,
so they agree everywhere. Thus $T'\to Z\to Y$ is
$G$-invariant. It descends along $T'/G=T$ to a morphism
$T\to Y$ extending $a$. This proves the second condition.
Any two such extensions agree on the dense open $T^\circ$;
normality of $T$ and separatedness of $Y$ give uniqueness.

For intrinsic uniqueness, suppose $(Y_1,U)$ and $(Y_2,U)$
satisfy the definition. Choose a covering pair $(Z_1,W_1)$
for the first. A finite normal integral $Z_1\to Y_1$ is the
normalization of $Y_1$ in its function field: on each affine,
its finite normal coordinate algebra equals the integral
closure in that field. Consequently the interior Galois action
extends to $Z_1$, and the invariant algebra is the normal
base algebra, so $Z_1/G_1=Y_1$. Its $K$-scheme is excellent and noetherian,
since it is finite over the finite-type $K$-scheme $Y_1$.
The second condition for $Y_2$ extends $W_1\to U\subset Y_2$
to $Z_1\to Y_2$. The extension is invariant under the finite
Galois group, by equality on $W_1$. Since $Y_1$ is normal,
the affine invariant-ring argument gives $Z_1/G_1=Y_1$,
so it descends to $Y_1\to Y_2$. Reverse the roles. The
two composites are identities on $U$, hence are identities
on the normal separated compactifications. The resulting
isomorphism is unique.

The same argument applied to $W_1\to U_1\to U_2$ extends
any map of interiors to $Y_1\to Y_2$. Its uniqueness makes
extension respect identities and composition. Applying this
to an arithmetic-group-induced interior isomorphism and its
inverse proves the last assertion. In particular, for the
strongly hyperbolic Artin neighbourhoods the needed interior
theorem is precisely Schmidt--Stix (2016), Corollary~1.6
\cite{A567_SchmidtStix2016}, already opened and applied above.
\end{proof}

\begin{corollary}[Known input, with proved compactification deduction]
\label{A:polycurve-four-dimensional-anabelian}
Let $p$ be a prime and let $K$ be a sub-$p$-adic field: that is,
$K$ embeds into a finitely generated extension of $\mathbb Q_p$.
Let $U,U'$ be geometrically connected hyperbolic polycurves over
$K$, and suppose that
$\min\{\dim U,\dim U'\}\leq4$.
Write $\Pi_U=\pi_1^{\mathrm{\acute et}}(U)$ and
$\Delta_U=\ker(\Pi_U\to G_K)$, and similarly for $U'$.
For their canonical normal projective compactifications from
Theorem~\ref{A:polycurve-coarse-compactification}, the natural map
\[
 \operatorname{Isom}_K\bigl((\overline U,U),
                         (\overline {U'},U')\bigr)
 \longrightarrow
 \operatorname{Isom}_{G_K}(\Pi_U,\Pi_{U'})/
                   \operatorname{Inn}(\Delta_{U'})
\]
is bijective. In particular every prescribed geometric-inner
class of an augmentation-preserving arithmetic group isomorphism
is realized by one and only one isomorphism of compactified pairs.
No chosen tower, split marking, level structure, saturation, or
strongly hyperbolic Artin hypothesis is imposed on these interiors.
\end{corollary}

\begin{proof}
Hoshi (2014), Theorem~B, pp.~156--157, proves exactly the bijection
\[
 \operatorname{Isom}_K(U,U')
  \xrightarrow{\;\sim\;}
 \operatorname{Isom}_{G_K}(\Pi_U,\Pi_{U'})/
                    \operatorname{Inn}(\Delta_{U'}).
\]
The more general Corollary~3.18, pp.~203--204, includes
hyperbolic orbi-polycurves; here we use only its scheme case
\cite{A567_Hoshi2014}. The definition of hyperbolic polycurve is
Definition~2.1, pp.~166--167: a sequence of smooth families of
hyperbolic curves, each obtained from a smooth proper curve by
removing a finite \emph{\'etale} divisor. It agrees with the
interior hypothesis of our compactification theorem.
Definition~3.1, p.~187, is precisely the field condition stated above.

For clarity, the full isomorphism assertion is not inferred just
from the abstract's object-reconstruction wording. In the proof of
Corollary~3.18, Claim~3.18.A, Hoshi starts with the \emph{specified}
isomorphism of arithmetic groups. After reversing it if necessary,
the smaller-dimensional target has dimension at most four.
Theorem~3.8 and Corollaries~3.12, 3.15 and 3.16, together with
Proposition~2.7, realize that homomorphism by
a quasi-finite dominant morphism and force the dimensions to agree.
The LFG hypothesis appearing in the dimension-three statement is
automatic for a hyperbolic polycurve by Proposition~2.7,
pp.~174--175. Apply the same argument to the inverse group
isomorphism. Proposition~3.2(ii), pp.~187--188, makes the two
compositions the identity, and also gives uniqueness. Thus the
input is the full bijection of isomorphism sets, with exactly the
geometric inner quotient displayed above.

Now restriction to the identified open interiors gives a bijection
\[
 \operatorname{Isom}_K\bigl((\overline U,U),
                         (\overline {U'},U')\bigr)
 \xrightarrow{\;\sim\;}
 \operatorname{Isom}_K(U,U').
\]
Surjectivity is the extension assertion of
Theorem~\ref{A:polycurve-coarse-compactification}; extend an interior
isomorphism and its inverse, and their compositions are the identity
because they are the identity on a dense open. The same dense-open
argument proves injectivity. Composing these two proved bijections
gives the result.

Every field finitely generated over $\mathbb Q$ satisfies the
sub-$p$-adic hypothesis for every prime $p$: choose a transcendence
basis, embed its rational-function field into
$\mathbb Q_p(t_1,\ldots,t_r)$, and extend the embedding over the
remaining finite algebraic field extension into an algebraic
closure. Its image lies in a finite extension of
$\mathbb Q_p(t_1,\ldots,t_r)$. This verifies the field hypothesis
in the arithmetic setting of 049(c).\qedhere
\end{proof}

The dimension condition here applies to arbitrary hyperbolic
polycurves. It is complementary to the unrestricted-dimensional
strongly hyperbolic Artin result used above. The statement preserves
the augmentation to a fixed $G_K$; it does not replace it by an
unaugmented profinite group, and it makes no assertion that every
outer homomorphism, as opposed to every outer isomorphism in the
displayed class, is geometric.

\paragraph{Exact scope of the answer to 049(c).}
The sheet says to \emph{define} an anabelian compactification,
inspired by stable moduli compactifications, and asks about
specialization. It does not quantify an assertion that every
smooth variety called anabelian has every kind of proper
boundary model canonically recoverable from its interior.
Definition~\ref{A:anabelian-compactification-definition} gives
an intrinsic category without naming a tower. Theorems
\ref{A:polycurve-coarse-compactification} and
\ref{A:intrinsic-anabelian-compactification} prove uniqueness,
functoriality, existence for \emph{all} characteristic-zero
hyperbolic polycurves, and arithmetic reconstruction in the
known strongly hyperbolic Artin class in every dimension.
Corollary~\ref{A:polycurve-four-dimensional-anabelian} further gives
full arithmetic isomorphism reconstruction for arbitrary hyperbolic
polycurves through dimension four over sub-$p$-adic fields. Theorem
\ref{A:compactification-obstruction} explains why ``any chosen
normal crossings compactification'' cannot replace this
definition. The specialization answer is
Theorem~\ref{A:log-specialization-higher}, with its precise
log-model and prime restrictions; an underlying singular
coarse fibre must not replace the log fibre.

The log-regular covering pair controls boundary ramification;
quotient stacks made from different covers may retain different
boundary stabilizers, so no cover-independent stack is asserted.
Existence of this intrinsic compactification for smooth varieties
outside the proved polycurve class is not asserted. That is a
possible additional research question, not a quantified claim
made in 049(c), and is not used to leave the stated definition
and comparison task unfinished.

\begin{proposition}[Proved: the smallest surface has a canonical log model]
\label{A:canonical-log-model-product}
For $U=(\mathbb P^1_K\setminus\{0,1,\infty\})^2$, let
$(Z,E)$ be any smooth projective compactification with reduced
simple normal crossings boundary. The graded ring
\[
 R(U)=\bigoplus_{m\ge0}H^0(Z,\omega_Z^{\otimes m}(mE))
\]
is independent of this choice, under the natural identification
as rational pluricanonical forms on $U$, and
\[
 \operatorname{Proj}R(U)\simeq\mathbb P^1_K\times\mathbb P^1_K.
\]
Thus the arithmetic interior group recovers this canonical log
model, although it does not recover an arbitrary chosen blowup.
\end{proposition}

\begin{proof}
We first check invariance under a single blowup $f:Z'\to Z$
at a boundary point. If the point lies on $r\in\{1,2\}$ local
boundary branches and $F$ is the exceptional curve, then
\[
 K_{Z'}+E'=f^*(K_Z+E)+(2-r)F.
\]
Indeed the Jacobian of $(x,u)\mapsto(x,xu)$ has factor $x$,
giving $K_{Z'}=f^*K_Z+F$; the reduced total boundary is
$E'=f^*E-(r-1)F$. The same calculation holds at a non-rational
closed point after an \'etale field extension and descends.
For every integer $a\ge0$,
$f_*\mathcal O_{Z'}(aF)=\mathcal O_Z$: a rational function with
possible poles only on $F$ is regular away from the centre on
$Z$, and extends across that codimension-two point because $Z$
is normal. The projection formula therefore identifies each
pluricanonical space above, compatibly with multiplication.

Two smooth projective surface compactifications have a common
resolution by point blowups; this is the opened Stacks Project,
Tags \texttt{0C5R}--\texttt{0C5S}, Lemmas~54.17.1--2. Because
the birational map is the identity on $U$, its indeterminacies
and all the required centres lie in the boundary. Point blowups
preserve the simple normal crossings property on a surface.
Applying the preceding calculation to the two sequences proves
the canonical equality of rings, not only equality of their
dimensions.

Now use $Z_0=\mathbb P^1\times\mathbb P^1$ with its six
coordinate boundary curves $E_0$. Then
$\omega_{Z_0}(E_0)=\mathcal O(1,1)$, so
\[
 R(U)=\bigoplus_{m\ge0}H^0(Z_0,\mathcal O(m,m)).
\]
Its four degree-one generators are $x_i y_j$ for $i,j\in\{0,1\}$.
They generate every bidegree-$(m,m)$ monomial. Their sole relation
is the Segre equation $z_{00}z_{11}-z_{01}z_{10}=0$:
the image is the rank-one matrix locus, and on each nonzero
coordinate chart this equation reconstructs the two projective
factors. Hence its Proj is the asserted surface. Functoriality of
rational pluricanonical forms under an isomorphism of $U$ gives
functoriality of this model. Applying the interior reconstruction
theorem proves the last assertion. The exact bidegree, blowup
discrepancy, Segre relation and sample Hilbert dimensions are
recorded in \texttt{code/A567\_compactification.py}; the argument
above proves invariance for all compactifications.
\end{proof}

\begin{theorem}[Known specialization theorem, with its prime restrictions]
\label{A:log-specialization-higher}
Let $R$ be a complete discrete valuation ring with fraction field
$F$, residue characteristic $p\ge0$, and its standard divisorial
log structure. Let $\mathcal X^{\log}\to(\operatorname{Spec}R)^{\log}$
be proper and log smooth, with fine saturated log structures and
geometrically connected generic fibre, and let $U$ be its
trivial-log locus. Write $s^{\log}$ for the induced log special
point. There is a Galois-compatible equivalence
\[
 \operatorname{Cov}_{\mathrm{\acute et}}(U_{\bar F})^{(p')}
 \simeq
 \operatorname{KCov}_{\mathrm{geom}}
       (\mathcal X_s^{\log}/s^{\log})^{(p')},
\]
and thus an outer isomorphism of the corresponding fundamental
groups. Superscript $(p')$ means finite quotients of order prime
to $p$; if $p=0$, all finite quotients are retained. The right
side uses log-geometric covers, obtained after all finite Kummer
extensions of the base log point, not ordinary covers of the
underlying singular special fibre.
\end{theorem}

\begin{proof}[Derivation from the opened primary theorem]
This is Kisin (2000), Theorem~1.4, actual p.~1103, with
Proposition~1.15 and Corollary~1.16, pp.~1115--1116; its
log-geometric notation is also set out in Lepage (2013),
Theorem~1.12, p.~1442. The first comparison extends a cover of
the trivial-log locus by normalization to a finite Kummer log
cover. Log purity supplies local root charts. Properness identifies
these covers with covers of the log special fibre by lifting and
algebraization. Passing through all tame extensions of $F$ removes
the base inertia and gives the log-geometric category. Kisin's
Proposition~1.15 is the needed assertion that prime-to-$p$
geometric covers descend over tame extensions, and his theorem
states Galois compatibility. Equivalences of pointed Galois
categories identify their automorphism groups; forgetting a
chosen point leaves the stated outer isomorphism. Lepage allows
any set of invertible primes, so when $p=0$ this includes every
prime. This derives the specified specialization law without
silently replacing a log special fibre by its underlying scheme.
\end{proof}

\paragraph{Comparison of the compactification statements.}
The saturated polystable statement retains its stronger conclusion
that the compactification itself is log regular. The quotient
construction removes saturation while retaining the original
interior, at the cost of allowing a normal coarse boundary.
The intrinsic toric extension property supplies a definition and
uniqueness criterion without retaining a chosen tower. Together
with the precise log specialization theorem, these answer the
definition and comparison asked for on 049(c). The boundary
blowup example refutes the different assertion that an arbitrary
chosen compactification is recoverable from the interior group.
No existence theorem for every possible broader class of
anabelian varieties, or unrestricted wild-prime specialization,
is being silently added to the manuscript's question.

% Original-source acquisition and scope comparison, 19 November 2026.
% Unnumbered replacement for the stale source-access paragraph in A.6.
\paragraph{The original 1990 and 1997 sources.}
The three original Crelle articles have now been opened from the
public G\"ottingen scans. Nakamura's volume~405 paper studies the
quotient with metabelian geometric kernel; Theorems~A--B retain
additional field or exceptional-unit hypotheses, while Theorem~C
recovers a multiplicative-difference invariant
\cite{A567_Nakamura1990Metabelian}.
His distinct volume~411 paper, Definition~3.3 and Theorem~3.4,
pp.~209--210, proves smooth cusp recognition over number fields using
self-normalizing procyclic subgroups, open arithmetic normalizer
image and the cyclotomic character. This confirms the original
provenance of the scoped smooth-cusp result used through Stix above.
Theorems~5.10 and~6.1, p.~215, reconstruct isomorphism of genus-zero
objects; Theorem~5.10 also realizes the induced cusp permutation.
They do not state that every prescribed group isomorphism is
induced by that curve isomorphism \cite{A567_Nakamura1990}.

Ihara--Nakamura's construction starts with a maximally degenerate
curve whose irreducible components are smooth projective lines.
Its canonical parametrization requires both a tangential structure
and an ordering of the nodes. For the chosen diagonal deformation,
Theorem~2$'$ compares kernels of Galois actions on tangential
pro-$\mathcal C$ fundamental groupoids.
Theorem~3B gives the universal pro-$\ell$ monodromy fixed-field
inclusion, and Corollary~4.2.2 gives equality for $n\geq1$ using
the previously known opposite inclusion
\cite[\S\S1.2--1.4, 3.6, 4.1--4.2, pp.~125--127,146--148]{A567_IharaNakamura1997}.
Immediately after Corollary~4.2.2, the article credits Takao with
the unpunctured case, without reproducing his proof.
Thus Oda's fixed-field problem is distinct from reconstruction in
049(b), as also explained by the opened Collas--Philip paper.
L6 identifies a relevant later programme; the tangential,
compactification and arithmetic-presentation assertions of
049(a--c) require the additional inputs and arguments supplied here.


\begin{theorem}[Proved: the two matrix presentations of 67/050]
\label{A:SL-GL-presentations}
The manuscript's presentations, with the involution relations made
explicit, are
\[
 \operatorname{SL}_2(\mathbb Z)
 =\langle\sigma,\rho\mid\sigma^2=\rho^3,\ \rho^6=1\rangle
\]
and
\[
 \operatorname{GL}_2(\mathbb Z)
 =\langle r_1,r_2,r_3\mid
 r_1^2=r_2^2=r_3^2=1,\
 (r_1r_2)^2=(r_2r_3)^3,\ (r_1r_2)^4=1\rangle.
\]
In the second presentation all three generators reverse orientation.
\end{theorem}

\begin{proof}
Use the matrices
\[
 S=\begin{pmatrix}0&-1\\1&0\end{pmatrix},\quad
 T=\begin{pmatrix}1&1\\0&1\end{pmatrix},\quad R=ST,
\]
and
\[
 r_1=\begin{pmatrix}-1&0\\0&1\end{pmatrix},\quad
 r_2=\begin{pmatrix}0&1\\1&0\end{pmatrix},\quad
 r_3=\begin{pmatrix}1&1\\0&-1\end{pmatrix}.
\]
They satisfy $S^2=R^3=-I$, $r_1r_2=S$, $r_2r_3=R$,
and all $r_i^2=I$, $\det r_i=-1$. These are checked with exact
integer arithmetic by \texttt{code/A567\_matrices.py}.

In the first abstractly presented group the common element
$c=\sigma^2=\rho^3$ is central and $c^2=1$. Quotienting by $c$
gives $C_2*C_3$. The opened source Milne (2017), Remark~2.14,
identifies this free product with $\operatorname{PSL}_2(\mathbb Z)$
under the images of $S,R$. Also $S,T$ generate
$\operatorname{SL}_2(\mathbb Z)$: left multiplication by their
powers implements the Euclidean algorithm on the first column of
any unimodular matrix, after which the remaining matrix is a power
of $T$, possibly times $-I=S^2$. Hence the homomorphism from the
abstract presentation to the matrix group is surjective, and its
kernel lies in $\langle c\rangle$. As $c$ maps to $-I\ne I$,
the kernel is trivial.

For the second group put $r=r_2$, $\sigma=r_1r_2$ and
$\rho=r_2r_3$. Its presentation is equivalently the first one
together with
\[
 r^2=1,\qquad r\sigma r=\sigma^{-1},\qquad
 r\rho r=\rho^{-1}.
\]
Conversely $r_1=\sigma r$ and $r_3=r\rho$ recover all the
original relations. Inversion of the two specified generators
preserves the first presentation and has order~2, so this is
$\operatorname{SL}_2(\mathbb Z)\rtimes C_2$. The displayed
$r_2$ implements that automorphism. The determinant map on
$\operatorname{GL}_2(\mathbb Z)$ has kernel
$\operatorname{SL}_2(\mathbb Z)$ and is split by $r_2$;
the presentation therefore gives the asserted matrix group.
\end{proof}

\begin{proposition}[Proved: the groups and exact sequences behind 67/051]
\label{A:four-puncture-MCG}
Let $V=(\mathbb Z/2)^2$ and
$\Gamma(2)=\ker(\operatorname{SL}_2(\mathbb Z)\to
\operatorname{SL}_2(\mathbb F_2))$. Then
\[
 \operatorname{Mod}_{0,4}\simeq
 V\rtimes\operatorname{PSL}_2(\mathbb Z),\qquad
 \operatorname{PMod}_{0,4}\simeq
 \Gamma(2)/\{\pm I\}\simeq F_2.
\]
The quotient of $\operatorname{Mod}_{0,4}$ by $V$, and also the
subgroup fixing one distinguished puncture, are each isomorphic to
$\operatorname{PSL}_2(\mathbb Z)$. The two exact rows are
\[
 1\longrightarrow F_2\longrightarrow\operatorname{Mod}_{0,4}
 \longrightarrow\mathfrak S_4\longrightarrow1,
\]
\[
 1\longrightarrow F_2\longrightarrow\operatorname{PSL}_2(\mathbb Z)
 \longrightarrow\mathfrak S_3\longrightarrow1.
\]
The map between their finite quotients is
$\mathfrak S_4\to\mathfrak S_4/V\simeq\mathfrak S_3$.
In particular $\operatorname{PMod}_{0,4}$ is not
$\operatorname{PSL}_2(\mathbb Z)$.
\end{proposition}

\begin{proof}
Farb and Margalit (2012), Proposition~2.7, computes the first
semidirect product using the double cover of the four-punctured
sphere by a torus branched over its four marked points. In that
construction $V$ consists of translations by the four two-torsion
points and the linear factor acts through reduction modulo~2.
Consequently the action on the four branch points is the affine
action on $\mathbb F_2^2$. Its image is
$\operatorname{AGL}_2(\mathbb F_2)$, which acts faithfully on
four elements and has order $4\cdot6=24$, hence is
$\mathfrak S_4$. The kernel consists exactly of zero translations
and matrices congruent to the identity modulo~2, modulo $\pm I$.
The stabilizer of the zero two-torsion point has zero translation
and is the full linear factor. These observations derive both
exact rows and the distinguished-puncture assertion from the
opened computation; no interpretation of the prime notation is
needed.

For completeness the kernel in the second row is free of rank~2.
The group $\operatorname{PSL}_2(\mathbb Z)=C_2*C_3$ acts on
its free-product tree. Reduction modulo~2 maps the two finite
factors injectively to subgroups of orders~2 and~3 in
$\mathfrak S_3$, and they generate that group. Its index-six
kernel therefore acts freely on the tree. The quotient graph has
$6/2+6/3=5$ vertices and $6$ edges, so its fundamental group is
free of rank $6-5+1=2$. A free group has no torsion; in contrast
$\operatorname{PSL}_2(\mathbb Z)$ contains elements of orders
2 and~3. They cannot be isomorphic. The matrix reduction and the
finite affine action are independently enumerated in
\texttt{code/A567\_matrices.py}.
\end{proof}

\begin{proposition}[Known, derived with its exceptions: the Birman sequence]
\label{A:Birman}
Let $S$ be an oriented surface of finite type with $\chi(S)<0$,
and let $p$ be a newly distinguished point in its interior. Then
\[
 1\longrightarrow\pi_1(S,p)\stackrel{\mathrm{Push}}\longrightarrow
 \operatorname{Mod}(S,p)\stackrel{\mathrm{Forget}}\longrightarrow
 \operatorname{Mod}(S)\longrightarrow1
\]
is exact. The new point must be preserved, even when the other
punctures may be permuted. The unrestricted assertion without
$\chi(S)<0$ is false; for the closed torus the forgetful map
$\operatorname{Mod}_{1,1}\to\operatorname{Mod}_{1,0}$ is an
isomorphism and the point-pushing map from $\mathbb Z^2$ is zero.
\end{proposition}

\begin{proof}
The first assertion is exactly Farb and Margalit (2012),
Theorem~4.6, applied with the prescribed new point. Its geometric
content is that an isotopy from a homeomorphism to the identity
traces a loop by the moving point; reversing that loop constructs
the inverse identification of the forgetful kernel with point
pushes. Injectivity follows from the induced conjugation action
on the centre-free group $\pi_1(S,p)$. The source proves that this
construction is independent of the isotopy. Restriction to the
pure groups gives the pure version on the page.

For a torus, a translation following any loop based at the origin
is an isotopy whose endpoint is the identity translation. Hence
every loop gives a trivial point push. Every torus mapping class
has a linear representative fixing the origin; its action on
homology identifies both the pointed and unpointed mapping class
groups with $\operatorname{SL}_2(\mathbb Z)$. Thus forgetting the
origin is an isomorphism. This verifies the exceptional case
directly and explains why the missing Euler-characteristic
condition cannot be suppressed.
\end{proof}

\paragraph{Pants and Lego presentations.}
The precise existing presentation theorem is stronger than
generation by elementary moves. Bakalov and Kirillov (2000),
Theorem~4.24, proves that their two-complex of markings of an
extended oriented surface is connected and simply connected, with
the $Z,F,B,S$ moves and the relations in \S\S4.7 and 4.22.
Their Theorem~5.1 gives the maximal-pants version with
$Z,F,A,B,S$. Boundary circles carry marked points in their
convention. To derive a groupoid presentation, take markings as
objects, oriented edges as generators, inverse edges as inverse
generators, and boundaries of the stated two-cells as relators.
Any morphism is represented by an edge path because the complex
is connected. Any two paths with the same ends differ by a loop,
and simple connectivity expresses that loop as a product of
boundaries of two-cells. Thus there are no further relations.
The passage to a mapping class group retains the group action and
its stabilizers. Here is its precise form. Fix a marking $M$ and
write $\Gamma=\operatorname{Mod}(\Sigma)$. Consider pairs
$(g,[p])$, where $g\in\Gamma$ and $p$ is an edge path from $M$
to $gM$, modulo the two-cell homotopies. Define
\[
 (g,[p])(h,[q])=(gh,[p*(gq)]).
\]
For each $g$ connectedness supplies $p$, and simple connectivity
makes its homotopy class unique. Thus $(g,[p])\mapsto g$ is a
group isomorphism. This is the automorphism group in the action
path groupoid, equivalently the fundamental group of the homotopy
quotient. It is not the fundamental group of $\mathcal M(\Sigma)$
itself, which is trivial. A presentation from quotient move types
must also include the stabilizers and their action relations;
discarding them is not justified. Bakalov--Kirillov's Remark~4.3
explicitly distinguishes a move from a surface homeomorphism,
and Theorem~3.7 proves that the stabilizer of a marking becomes
trivial after a specified boundary component is fixed.
Hatcher and
Thurston (1980), Theorem~1.1, supplies the analogous connected and
simply connected cut-system complex. These are precise topological
answers to the pants and Lego expectations on 050; their proofs do
not by themselves construct compatible arithmetic Galois actions.

\begin{proposition}[Proved reduction: orientation-reversing involutions]
\label{A:reflection-generators}
For an oriented closed surface with finitely many marked points,
the extended mapping class group has a finite generating set
consisting of orientation-reversing involutions. It therefore has
a finite presentation on such generators. This assertion concerns
marked points; it is not asserted for a group required to fix a
nonempty boundary pointwise.
\end{proposition}

\begin{proof}
The finite twist generating theorem for pure mapping class groups,
and the addition of finitely many half twists for permutations of
the marked points, are proved in Farb and Margalit (2012),
Theorems~4.9, 4.11 and \S4.4.4. Let $x_1,\ldots,x_s$ be such
generators. Each twist curve or two-point half-twist disc has a
symmetric surface model: cut along the curve, realize each
resulting oriented surface as a sphere with symmetrically attached
handles, and put its remaining marked points on the reflection
arcs. Gluing the reflected boundary circles gives an
orientation-reversing reflection $a_i$ preserving the twist curve.
For a half twist use a disc whose two marks lie on a reflection
axis and extend its reflection to a symmetric model of the
complement. The classification of oriented surfaces with marked
points transfers each model to the specified surface. Thus
$a_i^2=1$ and $a_ix_ia_i=x_i^{-1}$. The element
$b_i=x_i a_i$ is also an orientation-reversing involution, since
$b_i^2=x_i(a_ix_ia_i)=1$, and $x_i=b_i a_i$.

Fix one additional reflection $r$. Products of the $a_i,b_i$
generate the orientation-preserving group, and adjoining $r$
gives its full extension by orientation reversal. This is a finite
generating set of the required kind. Finite presentability of
mapping class groups is proved by the opened pants/marking
complex, or Farb--Margalit \S5.3. A finite-index extension is
finitely presented: add the generator $r$, the relation $r^2=1$,
and the finitely many conjugation relations for the old generators.
Finally make the above substitution explicit as Tietze moves:
adjoin symbols $a_i$ with relations identifying them with their
words in the old generators and $r$, adjoin $b_i=x_i a_i$, add
the consequent relations $a_i^2=b_i^2=1$, and eliminate
$x_i=b_i a_i$. This is a finite presentation entirely in the
displayed types of involution. It asserts existence, not an
unproved uniform short list of all its relators. Genus one has
the completely explicit three-reflection presentation of
Theorem~\ref{A:SL-GL-presentations}.
\end{proof}

\paragraph{The arithmetic meaning of the pants and Lego presentations.}
The comparison in the preceding paragraph can be made arithmetic.
The relevant objects are profinite fundamental \emph{groupoids} of
$\mathcal M_{g,n}$ based at rational tangential points along maximally
degenerate curves. A quilt is a pants decomposition with three seams on
each pair of pants; it specifies the extra tangential choices. Changes
of quilt are retained as paths between objects. This distinction is
essential when comparing a local geometric move with a based group
automorphism.

\begin{theorem}[Known, with explicit derivation: arithmetic quilt moves]
\label{A:arithmetic-quilts}
For every stable type $2g-2+n>0$, the profinite Teichm\"uller groupoids
with the above tangential objects admit the natural $G_{\mathbb Q}$
action in pants/quilt coordinates. Fusing moves are controlled by the
action on $\mathcal M_{0,4}$, simple moves by the elliptic moduli problem
with a nonzero tangent vector, and local Dehn twists and half-twists by
the cyclotomic character. The relations are the usual topological
quilt relations, with changes of base object and stabilizers retained.
At a fixed quilt this gives a lift to automorphisms of the profinite
pure mapping class group; at each Galois element a different
quilt changes that automorphism by a specified inner factor.
The factor may depend on the Galois element: the assertion is
equality of the outer actions, not conjugacy of the two entire
representations by one fixed inner automorphism. The resulting
outer action is compatible with
inclusion of quilted subsurfaces and with filling boundary circles by
marked discs.
\end{theorem}

\begin{proof}
We separate the topological presentation from its arithmetic realization.
Nakamura and Schneps (2000), Definition~1.1 and Theorems~1.2--1.4
\cite{A567_NakamuraSchneps2000},
construct a group $\mathrm I\Gamma\subset\widehat{\mathrm{GT}}$ containing
$G_{\mathbb Q}$ and, for each quilt $Q/P$ on a surface $\Sigma$, a
homomorphism
\[
 \rho^\Sigma_{Q/P}:\mathrm I\Gamma\longrightarrow
       \operatorname{Aut}\widehat\Gamma(\Sigma).
\]
Their \S7 defines the pro-words associated to $A$- and $S$-moves;
\S8 proves independence of the move chain; \S9 verifies the mapping
class relations. Thus the action exists on the presented group, and
not only on a free group on twist symbols. Equations~(10.1)--(10.4)
give its change under a different quilt explicitly as
$\rho_{Q'}(\sigma)=\operatorname{Inn}(c_\sigma)\rho_Q(\sigma)$.
Theorem~1.4
and its proof in \S10 give compatibility with subsurface inclusion.
The first paragraph of \S10 also proves the filling assertion: the
kernel of filling marked discs is generated by the corresponding
boundary twists, each sent to its $\lambda$th power. Restrict
$\mathrm I\Gamma$ to $G_{\mathbb Q}$, where $\lambda=\chi(\sigma)$.

For the identification with the actual arithmetic action, we use
Ichikawa's explicit geometric realization, rather than infer it from
the abstract group construction. His 2002 author proceedings paper,
Theorem~2 and \S3, pp.~55--57, constructs the real orbifold of local
moves and its integral tangential parameters. Theorems~3--4 and
\S4.3, pp.~58--59, compute the arithmetic action. If $q$ is a fusing
path and $p$ the standard path from the tangential point at $0$ to the
one at $1$ on $\mathbb P^1\setminus\{0,1,\infty\}$, put
$f_\sigma=\sigma(p)p^{-1}$. In the author's path convention,
\[
 \sigma(q)q^{-1}=f_\sigma(\delta_0,\delta_1),
\]
where $\delta_0,\delta_1$ are the two local Dehn twists. For a simple
move $s=\iota(r)$, with $r$ the corresponding path in the tangent
version of $\mathcal M_{1,1}$, his Theorem~4 gives
$\sigma(s)=\iota(\sigma(r))$. The tangent version retains the central
boundary twist; suppressing it would confuse the one-holed torus group
with $\operatorname{SL}_2(\mathbb Z)$. Finally \S4.3 gives
$\sigma(\delta)=\delta^{\chi(\sigma)}$ for a local twist or half-twist.
These exact formulas are restated in Ichikawa (2006), Theorems~2--3,
pp.~3--4 \cite{A567_Ichikawa2002,A567_Ichikawa2006}.
Both author texts were opened, including the construction
and the theorem statements; the 2002 text expressly reviews the
results of the later 2003 paper without reproducing all its proofs.

The move paths generate the topological groupoid by the already
opened pants/marking theorem. Their profinite completion topologically generates
the corresponding profinite groupoid, and the displayed arithmetic
formulas determine a continuous functor on it: they are images of the
actual $G_{\mathbb Q}$ action on the geometric fundamental groupoid.
Consequently they preserve every relation and every stabilizer.
Restricting to loops at one object gives the based action above;
changing the object gives precisely the stated inner factor. This
derives the arithmetic interpretation, with the extra tangential
data specified, from the cited theorems.
\end{proof}

\paragraph{Involutions require a real-place interpretation.}
An orientation-reversing involution belongs to an extended mapping
class group. Over a chosen real closure, complex conjugation supplies
such an arithmetic extension. It is not an additional geometric
generator of the complex moduli stack's fundamental group. In
particular, the following tempting stronger interpretation of the
three presentations is false already in type $(0,4)$.

\begin{proposition}[Proved: the exact real-place obstruction]
\label{A:real-Galois-obstruction}
Let $U=\mathcal M_{0,4}=\mathbb P^1_{\mathbb Q}\setminus
\{0,1,\infty\}$, choose the rational base point $2$, and set
\[
 \Delta=\pi_1(U_{\overline{\mathbb Q}},2)\simeq\widehat F_2.
\]
Fix an embedding $\overline{\mathbb Q}\hookrightarrow\mathbb C$,
write $c$ for its complex conjugation, and let
$H_c=\Delta\rtimes\langle c\rangle$, using the action supplied by the
base point. This is the arithmetic fundamental group over the chosen
real closure, and is the profinite completion of the corresponding
extended pure four-point mapping class group. An element
$\sigma\in G_{\mathbb Q}$ has its natural outer action on $\Delta$
induced by an automorphism of $H_c$ preserving $\Delta$ if and only if
$\sigma\in\{1,c\}$. Thus the whole $G_{\mathbb Q}$ action cannot be
extended to this single fixed real extended group with the indicated
compatibility.

There is instead a canonical outer action on the family
$\{H_{\sigma c\sigma^{-1}}\}_{\sigma\in G_{\mathbb Q}}$ by
isomorphisms between its members. With the chosen rational base point
this action can be made strict inside
$\Delta\rtimes G_{\mathbb Q}$.
\end{proposition}

\begin{proof}
First the outer representation
$\rho:G_{\mathbb Q}\to\operatorname{Out}\Delta$ is faithful. Here is
the derivation needed, to avoid assuming faithfulness of an inner
lift is the same assertion. For $\sigma\ne1$, choose an algebraic
number $j$ with $\sigma(j)\ne j$, and an elliptic curve $E$ over
$\overline{\mathbb Q}$ having invariant $j$. Such a curve exists; for
$j\ne0,1728$ one may take
\[
 E:\ y^2=x^3-3kx-2k,\qquad k=\frac{j}{j-1728},
\]
whose discriminant is nonzero and whose invariant is $j$.
The chosen $j$ is automatically different from the rational numbers
$0$ and $1728$. Belyi's theorem, in the fully opened proof of
K\"ock (2004), Theorem~3.3 and Lemmas~3.4--3.6
\cite{A567_Kock2004}, gives a finite map
$E\to\mathbb P^1$ ramified only at $0,1,\infty$. Removing those fibres
gives a connected finite \emph{\'etale} cover of $U_{\overline{\mathbb Q}}$.
If $\rho(\sigma)=1$, its conjugacy class of open subgroups, hence its
cover isomorphism class, is fixed. An isomorphism with the conjugate
cover extends to their smooth projective compactifications. It would
give $E\simeq E^\sigma$, contradicting $j\ne\sigma(j)$. Thus $\rho$
is injective.

Put $r=\rho(c)$, represented by the actual base-point involution of
$\Delta$. Suppose $A$ is an automorphism of $H_c$ preserving $\Delta$,
and its restriction has outer class $\rho(\sigma)$. Since
$H_c/\Delta=C_2$ has trivial automorphism group, $A(c)=d c$ for some
$d\in\Delta$. Applying $A$ to the conjugation relation gives
\[
 A|_\Delta\, r\,(A|_\Delta)^{-1}=\operatorname{Inn}(d)r.
\]
In $\operatorname{Out}\Delta$ this says
$\rho(\sigma c\sigma^{-1})=\rho(c)$. Faithfulness implies
$\sigma c\sigma^{-1}=c$.

The centralizer of $c$ in $G_{\mathbb Q}$ is exactly $\{1,c\}$.
Indeed an element commuting with $c$ preserves the real algebraic
numbers $\mathbb Q^{\mathrm{rc}}=\overline{\mathbb Q}\cap\mathbb R$.
This real closed field has its unique ordering, since its positive
elements are precisely its nonzero squares. Every field automorphism
preserves that ordering and fixes $\mathbb Q$, hence fixes every real
algebraic number by its rational cut. It therefore lies in
$\operatorname{Gal}(\overline{\mathbb Q}/\mathbb Q^{\mathrm{rc}})
=\{1,c\}$. Conversely, the identity and conjugation by $c$ on $H_c$
induce the required two outer actions. This proves the equivalence.

The topological identification used in the statement can also be
seen directly. Real conjugation on the four-point moduli curve fixes
the point $2$ and lifts to an involution of its contractible
Teichm\"uller space. Adding this lift to the deck group $F_2$ gives
the extended pure mapping class group $F_2\rtimes C_2$. Finite
\'etale covers over the real closure are precisely complex finite
covers with this conjugation descent, so their group is
$\widehat F_2\rtimes C_2$. Completion preserves this semidirect
product because the first factor has finite index.

For a finite witness that $c$ is not central, let
$L=\mathbb Q(\alpha,\zeta_3)$ with $\alpha^3=2$ real. Eisenstein's
criterion and $\zeta_3\notin\mathbb Q(\alpha)\subset\mathbb R$ give
$[L:\mathbb Q]=6$. It is the splitting field of $x^3-2$, so its
Galois group is $\mathfrak S_3$. On its roots
$(\alpha,\alpha\zeta_3,\alpha\zeta_3^2)$, conjugation is $(23)$,
whereas the automorphism $\alpha\mapsto\alpha\zeta_3$,
$\zeta_3\mapsto\zeta_3$ is $(123)$. They do not commute. The exact
permutation certificate is
\texttt{code/A567\_real\_Galois.py}, with output in
\texttt{logs/A567\_real\_Galois.json}.

Finally use the rational base point to write
$\pi_1(U_{\mathbb Q},2)=\Delta\rtimes G_{\mathbb Q}$. Conjugation by
$(1,\sigma)$ carries $H_c$ to $H_{\sigma c\sigma^{-1}}$ and induces
the original $\sigma$-action on $\Delta$. These maps compose exactly.
Without a splitting, different lifts of $\sigma$ change them by
inner automorphisms of the target, giving the canonical outer
groupoid formulation. This proves the replacement, not merely the
failure of an action on a fixed $H_c$.
\end{proof}

\paragraph{Scope of the comparison with sheet 050.}
The phrases ``pants'', ``Lego'' and ``involutions'' on the sheet are
expectations about types of presentation, not definitions of one
common arithmetic category. The topological results above establish
all three types. Theorem~\ref{A:arithmetic-quilts} supplies the
arithmetic groupoid interpretation of pants and Lego; the complex
conjugation extensions just proved give a real-place meaning to the
involutions. Proposition~\ref{A:real-Galois-obstruction} rules out
putting these three interpretations into a single fixed extended
group with its unchanged full $G_{\mathbb Q}$ action. No identification
of $G_{\mathbb Q}$ with the entire automorphism group of a combinatorial
tower is needed or claimed.

\paragraph{Status of A.6.}
The matrix presentations and corrected discrete mapping class
sequences are proved, including the necessary exceptional case and
the prime/exclamation reading. The literal procyclic criterion
and literal three-conjugacy-class reading are refuted, with precise
replacements. Logarithmic graphicity is derived from opened primary
theorems. The tangential interpretation of $[Y,N]$ and its full isomorphism
reconstruction have been established with the local comparison and
Kummer steps explicit. The pants and Lego presentations have the arithmetic groupoid
interpretation proved above; involutions have the real-place
interpretation, with the precise obstruction to an unchanged
full Galois action on one fixed real extended group. This answers
A.6(b)'s requested comparison with known presentations, without
requiring an identification of the entire automorphism group of
a combinatorial tower with the absolute Galois group. For 049(c),
the compactification and specialization statements above retain
their exact categories and hypotheses. The sheet asks for a
definition and its specialization behavior; it does not itself
assert that every smooth anabelian variety has an arbitrary chosen
boundary model recoverable from its interior group.

\section{A.7. Real elliptic curves and integral involutions}

\paragraph{Meaning of the parameter.}
We use $a$ for a ratio of squared lengths, exactly as in 67/097--099.
Thus $a>0$ and the associated period ratios are
$i\sqrt a$ and $(1+i\sqrt a)/2$. We use the usual modular invariant
$j(i)=1728$. The complex uniformization and the correspondence of
homomorphisms with complex-linear maps preserving the lattices are
the results of Milne (2021), Chapter~III, \S3, especially
Proposition~3.3 and Theorem~3.10; the source was opened.
The elementary classification below is proved directly from that
correspondence. Alling's 1981 book, \S0.10 and the contents of
Chapter~12, were opened and identify the same annulus, M\"obius-band
and Klein-bottle classification. The full 1971 Alling--Greenleaf
chapter was not accessible in this run; we do not pretend to have
checked a theorem inside it.

\begin{theorem}[Proved: classification of the real elliptic groups]
\label{A:real-elliptic}
Every elliptic curve over $\mathbb R$, with its origin, is isomorphic
over $\mathbb R$ to exactly one member of the following two disjoint
families, with complex conjugation as real structure:
\[
 E_a(\mathbb C)=\mathbb C/(\mathbb Z+i\sqrt a\,\mathbb Z),
 \qquad
 E'_a(\mathbb C)=\mathbb C/
       \bigl(\mathbb Z+\tfrac{1+i\sqrt a}{2}\mathbb Z\bigr),
 \qquad a>0.
\]
For a fixed complex elliptic curve admitting a real structure, there
are two real forms. Except for the square lattice, they are either
$\{E_a,E_{1/a}\}$ or $\{E'_a,E'_{1/a}\}$ with $a>1$.
For the square lattice they are $E_1$ and $E'_1$.
For the hexagonal lattice they are $E'_3$ and $E'_{1/3}$.
The corresponding modular functions are
\[
 j(E_a)=j(i\sqrt a),\qquad
 j(E'_a)=j\bigl((1+i\sqrt a)/2\bigr).
\]
In each family $j(a)=j(1/a)$. In particular the guess
$j=c\log a$ on 67/097 is false for every constant $c$.
\end{theorem}

\begin{proof}
Uniformize a real elliptic curve by $\mathbb C/\Lambda$. Its
origin-preserving antiholomorphic involution lifts to
$z\mapsto u\overline z$ with $u\overline u=1$. Choose $c\ne0$ such
that $cu=\overline c$; multiplication by $c$ makes the involution
ordinary conjugation. Thus $\Lambda$ is invariant under conjugation.
Put $P=\Lambda\cap\mathbb R$ and $Q=\Lambda\cap i\mathbb R$.
They have rank one: the real-linear involution has one-dimensional
positive and negative eigenspaces and its integral matrix supplies
rational eigenvectors. Moreover
\[
 2\Lambda\subset P\oplus Q\subset\Lambda.
\]
The subgroup $\Lambda/(P\oplus Q)$ of
$\tfrac12(P\oplus Q)/(P\oplus Q)\simeq\mathbb F_2^2$ meets neither
coordinate axis, because $P,Q$ already contain all lattice points on
their respective axes. It is therefore either zero or the diagonal
line. After a positive real rescaling, primitive generators are
$1$ and $i\sqrt a$; the zero and diagonal cases give the two
displayed lattices.

An origin-preserving isomorphism between two such real elliptic
curves lifts to $z\mapsto dz$. Commutation with conjugation on the
torus implies $d=\overline d$: the continuous difference
$(d-\overline d)\overline z$ takes values in a discrete lattice and
vanishes at zero. Thus $d$ is real and rescales both primitive
eigenvectors by the same absolute factor. Their squared-length
ratio and the index of their sum in the lattice are unchanged.
This proves uniqueness, including disjointness of the families.

For completeness we determine all forms of a fixed complex curve.
The finite group of complex scalars preserving $\Lambda$ is cyclic
of order $n=2,4,$ or $6$. Indeed such a scalar has modulus~1,
and its trace on the integral rank-two lattice is
$u+\overline u\in\mathbb Z\cap[-2,2]$; the possibilities give
precisely these cyclic groups. Once one antiholomorphic involution
$\sigma$ is chosen, all the others are $u\sigma$. Since
$\sigma u\sigma^{-1}=\overline u=u^{-1}$, all these maps square to
one, and conjugation by a holomorphic scalar $v$ replaces $u$ by
$v^2u$. Hence there are exactly two conjugacy classes.
For $n=2$ or $6$, $\sigma$ and $-\sigma$ represent the two classes.
Their invariant and anti-invariant primitive vectors are exchanged,
so $a$ becomes $1/a$ and the lattice index is unchanged. If $a=1$
the lattice is square, so this case has already been excluded.
For $n=4$, $\sigma$ and $-\sigma$ are conjugate; taking
$\Lambda=\mathbb Z+i\mathbb Z$, conjugation is split and
$i\sigma$ has primitive eigenvectors $1+i$ and $1-i$, whose sum
has index two. Their squared lengths are equal. The two forms are
therefore $E_1$ and $E'_1$.

For the hexagonal lattice choose unit vectors $u,v$ with angle
$60$ degrees, and reflect in the line $\mathbb Ru$. The primitive
anti-invariant vector is $2v-u$, the eigenlattice has index two,
and
\[
 |2v-u|^2=4|v|^2-4\operatorname{Re}(v\overline u)+|u|^2
         =4-2+1=3.
\]
The parameter is $3$, not $\sqrt3$. The opposite reflection gives
$1/3$. This proves the correction at all three occurrences on
67/098. The displayed $j$ formulas follow from the period ratios.
The two real forms of a given complex curve have the same $j$, so
$j(a)=j(1/a)$. If $j=c\log a$, this equality would give
$2c\log a=0$ for all $a>0$, forcing $c=0$. But $j(1)=1728$ in
both families, so the identically zero possibility also fails.
\end{proof}

\begin{proposition}[Proved: the count on 67/100 with its proviso]
\label{A:pid-involutions}
Let $k$ be a principal ideal domain of characteristic different
from~2 such that $k_0=k/2k$ is a field. Conjugacy classes of
$\sigma\in\operatorname{GL}_2(k)$ with $\sigma^2=1$ and
$\sigma\ne\pm1$ consist of one split class and a collection of
nonsplit classes indexed by
\[
 k_0^\times/\operatorname{im}(k^\times\longrightarrow k_0^\times).
\]
There are exactly two classes if and only if the unit map is
surjective. For $k=\mathbb Z[(1+\sqrt{-11})/2]$ there are four
classes in total: one split and three nonsplit. Thus the struck
unit-surjectivity proviso on 67/100 is necessary.
\end{proposition}

\begin{proof}
Let $P=\ker(1-\sigma)$ and $Q=\ker(1+\sigma)$ inside $M=k^2$.
Over the fraction field the eigenvalues are $1,-1$, each once,
since otherwise $\sigma=\pm1$. As torsion-free rank-one submodules
of a finite free module over a PID, $P$ and $Q$ are free of rank one.
The inclusion $2M\subset P\oplus Q\subset M$ gives a subspace
\[
 m=2M/(2P\oplus2Q)\subset P/2P\oplus Q/2Q
\]
meeting neither coordinate line. Conversely the inverse image
$M'$ of any such $m$ inside $P\oplus Q$, with the involution
$(p,q)\mapsto(p,-q)$, reconstructs $(M,\sigma)$ using the
isomorphism $M\xrightarrow{2}M'$. Its positive and negative
eigenspaces are $2P$ and $2Q$, canonically identified with
$P,Q$ by multiplication by~2. This also proves the correspondence
on morphisms: a commuting map preserves the two eigenspaces, and
a pair of maps preserving $m$ induces a map of the reconstructed
modules. Thus no additional conjugacies have been discarded.

The only possibilities for $m$ are zero and a graph line of
nonzero slope in $k_0^2$. Changing generators of $P,Q$ multiplies
that slope by the reduction of a unit of $k$; every such change is
available. This gives exactly the stated quotient. Matrices
$\bigl(\begin{smallmatrix}1&t\\0&-1\end{smallmatrix}\bigr)$,
where $t$ runs through lifts of the nonzero residue classes,
represent the nonsplit classes (the slope may be replaced by its
inverse without changing this parametrization).

Put $\omega=(1+\sqrt{-11})/2$. The norm on $k$ is
$N(a+b\omega)=a^2+ab+3b^2$. This ring is norm-Euclidean:
for any complex number, first choose the closest multiple of
$\operatorname{Im}\omega=\sqrt{11}/2$ for its imaginary part,
then the closest integer for the residual real part. The squared
distance from the chosen lattice point is at most
$11/16+1/4=15/16<1$. Applied to a quotient of two elements of the
fraction field, this is the Euclidean division inequality, so $k$
is a PID. Since
\[
 N(a+b\omega)=(a+b/2)^2+11b^2/4,
\]
the units are exactly $\pm1$. Reduction modulo~2 gives
$k_0=\mathbb F_2[\bar\omega]/(\bar\omega^2+\bar\omega+1)
\simeq\mathbb F_4$. Both units reduce to~1; its three nonzero
elements therefore yield three different nonsplit classes.
\end{proof}

\begin{proposition}[Refuted: connectedness alone does not repair 67/091]
\label{A:eigenmodule-correction}
The attached correction R5, which adds connectedness of
$\operatorname{Spec}k$ to the rank-two projective corollary on
67/091, is insufficient. Even for a connected commutative
2-torsion-free ring, the positive and negative eigenspaces of a
nontrivial involution on a free rank-two module need not be
invertible modules. The corollary is valid over a PID, as in
Proposition~\ref{A:pid-involutions}; it is also valid over a
connected ring with $2$ invertible, provided the involution is
neither $1$ nor $-1$.
\end{proposition}

\begin{proof}
Take $k=\mathbb Z[t]/(t^2)$ and
$\sigma=\bigl(\begin{smallmatrix}1&t\\0&-1\end{smallmatrix}\bigr)$
on $M=k^2$. The ring is commutative, multiplication by~2 is
injective, and its spectrum is connected because its reduction
is $\mathbb Z$. Direct multiplication gives $\sigma^2=1$,
and $\sigma\ne\pm1$. The positive eigenspace is $ke_1$.
The negative eigenspace consists of pairs satisfying
$2x+ty=0$. Projection to $y$ identifies it with the ideal
$I=(2,t)$: writing $y=a+bt$, the equation requires $a$ to be even,
and then uniquely $x=-(a/2)t$. At the maximal ideal
$\mathfrak m=(2,t)$, the vector space
$I/\mathfrak m I= (2,t)/(4,2t)$ over $\mathbb F_2$ has independent
generators represented by $2,t$. Hence $I_{\mathfrak m}$ requires
two generators and is not a free rank-one module. It is therefore
not invertible.
The exact dual-number matrix identity and the four normal forms of
$(2,t)/(4,2t)$ are checked in \texttt{code/A567\_involutions.py}.

For the corrected assertion with $2$ invertible, the idempotents
$(1+\sigma)/2$ and $(1-\sigma)/2$ split $M$ into projective
eigenspaces. Their ranks are locally constant and hence constant
on a connected spectrum. Neither rank can be zero, since that
would force $\sigma=\mp1$ everywhere. Their ranks sum to two,
so each is one. The PID case was proved above.
\end{proof}

\begin{theorem}[Proved: topology and torsors, 67/099--101]
\label{A:real-topology}
For every $a>0$,
\[
 E_a(\mathbb R)\simeq (\mathbb R/\mathbb Z)\times\mathbb Z/2\mathbb Z,
 \qquad E'_a(\mathbb R)\simeq\mathbb R/\mathbb Z,
\]
as compact real Lie groups. Moreover
\[
 H^1(\mathbb R,E_a)\simeq\mathbb Z/2\mathbb Z,
 \qquad H^1(\mathbb R,E'_a)=0.
\]
The quotient Klein surfaces of the two elliptic groups are
respectively an annulus and a M\"obius band. The unique nonzero
torsor under $E_a$ has no real points; its complex curve modulo
its real structure is a Klein bottle, a surface without boundary.
Thus the deletion of ``sans'' on 67/101 must be reversed.
\end{theorem}

\begin{proof}
Write $b=\sqrt a$. In the split lattice
$\Lambda=\mathbb Z+ib\mathbb Z$, a point $z=x+iby$ is real
modulo $\Lambda$ exactly when $2y\in\mathbb Z$. This gives the
two circles $y=0,1/2$ and the displayed direct product, with
$(ib/2)$ as an order-two point on the second circle. In the
nonsplit lattice put $u=1,v=(1+ib)/2$. Conjugation sends
$u$ to $u$ and $v$ to $u-v$. The point $xu+yv$ is fixed modulo
the lattice exactly when $y\in\mathbb Z$, so every fixed class
is represented on the real axis and there is one circle.

Let $\gamma=\operatorname{Gal}(\mathbb C/\mathbb R)$ and let
$\sigma$ be its nontrivial element. The exact sequence of
$\gamma$-modules
$0\to\Lambda\to\mathbb C\to E(\mathbb C)\to0$ gives
\[
 H^1(\gamma,E(\mathbb C))\simeq H^2(\gamma,\Lambda)
 =\Lambda^\gamma/(1+\sigma)\Lambda.
\]
Here the higher cohomology of the real vector space $\mathbb C$
vanishes by averaging over the two-element group; the last
formula follows from its two-periodic resolution, whose
differentials alternate $1-\sigma$ and $1+\sigma$. In the
split case the denominator is $2\mathbb Z$ inside
$\Lambda^\gamma=\mathbb Z$; in the nonsplit case
$(1+\sigma)v=u$, so the denominator is all of
$\Lambda^\gamma$. This proves both cohomology assertions.
The standard descent classification identifies this $H^1$ with
isomorphism classes of $E$-torsors: a cocycle $c$ changes
conjugation to $z\mapsto\overline z+c$, and changing the
complex trivialization by a translation changes $c$ by a
coboundary. This gives a direct verification of that
classification in the present case.

In the split case the nonzero class is represented by $c=1/2$.
The map $z\mapsto\overline z+1/2$ squares to a lattice
translation and has no fixed point modulo $\Lambda$, since
the real part of $z-\overline z$ is zero whereas that of
$1/2+m+inb$ cannot be zero for an integer $m$.
Reflection of the rectangular torus in its real axis gives an
annulus. The nonsplit reflection has one fixed circle and
orientation-reversing quotient, yielding a M\"obius band:
one may use the band between adjacent lifts of that circle as
a fundamental domain, whose two remaining sides are glued
with a reversal. The free glide reflection in the split case
identifies the ends of a cylinder with a reversal, giving a
Klein bottle. These explicit fundamental polygons establish
the asserted topologies and the absence of boundary in the
torsor case.
\end{proof}
